% ---------------------------------------------------------------------------
% Draft. The camera-ready must use the official ICLR 2028 style file
% (iclr2028_conference.sty); this file uses a neutral article layout so the
% text, tables and numbers can be reviewed before the style is applied.
% ---------------------------------------------------------------------------
\begin{filecontents*}[overwrite]{figures/data/fig1-ratchet-floor.dat}
% fig1-ratchet-floor.dat - generated by data/make_figure_data.py
K meas8 meas11 Spert Sterm pred8ratchet pred8term pred11ratchet pred11term
5 0.393675 0.138809 609.195 609.195 0.337913 0.337913 0.0422391 0.0422391
6 0.263191 0.0339435 609.615 60.4279 0.338146 0.0335186 0.0422683 0.00418983
7 0.257591 0.0336181 609.615 1 0.338146 0.000554687 0.0422683 6.93359e-05
\end{filecontents*}

\begin{filecontents*}[overwrite]{figures/data/fig1-ratchet-prefix.dat}
% fig1-ratchet-prefix.dat - generated by data/make_figure_data.py
t S
1 22.1669
2 53.2531
3 173.257
4 417.772
5 609.615
6 69.0011
7 0.00248899
\end{filecontents*}

\begin{filecontents*}[overwrite]{figures/data/fig2-delta.dat}
% fig2-delta.dat - generated by data/make_figure_data.py
safety delta4 delta5 delta8
1 0 0 0
1.01 0.0075 0.01 0.01
1.02 0.0175 0.02 0.02
1.035 0.0475 0.035 0.0375
1.05 0.0625 0.0575 0.055
\end{filecontents*}

\begin{filecontents*}[overwrite]{figures/data/fig2-edge.dat}
% fig2-edge.dat - generated by data/make_figure_data.py
step edge
1 1
2 1.92062
3 1.71472
4 1.28578
5 1.01793
6 1.00004
\end{filecontents*}

\begin{filecontents*}[overwrite]{figures/data/fig2-escape.dat}
% fig2-escape.dat - generated by data/make_figure_data.py
step exact e5m2
0 0.704775 0.709424
1 1.87179 1.92972
2 1.64572 1.96039
3 1.18577 3.17385
4 1.01026 78.3153
5 0.999997 1.04778e+09
6 1 4.91673e+44
\end{filecontents*}

\begin{filecontents*}[overwrite]{figures/data/fig3-bf16.dat}
% fig3-bf16.dat - generated by data/make_figure_data.py
bits err
24 0.716996
32 0.296493
40 0.0243381
48 0.0160026
56 0.0161829
\end{filecontents*}

\begin{filecontents*}[overwrite]{figures/data/fig3-e4m3.dat}
% fig3-e4m3.dat - generated by data/make_figure_data.py
bits err
15 0.722627
20 0.313055
25 0.129298
30 0.124283
35 0.125293
\end{filecontents*}

\begin{filecontents*}[overwrite]{figures/data/fig3-e5m2.dat}
% fig3-e5m2.dat - generated by data/make_figure_data.py
bits err
12 0.734967
16 0.363594
20 0.248821
24 0.240794
28 0.238433
\end{filecontents*}

\begin{filecontents*}[overwrite]{figures/data/fig3-frontier.dat}
% fig3-frontier.dat - generated by data/make_figure_data.py
bits err
12 0.734967
15 0.722627
16 0.363594
20 0.248821
24 0.240794
25 0.129298
30 0.124283
38 0.121183
40 0.0243381
48 0.0160026
\end{filecontents*}

\begin{filecontents*}[overwrite]{figures/data/fig3-ladder.dat}
% fig3-ladder.dat - generated by data/make_figure_data.py
bits err body six_m m
32 0.240194 e5m2 0 6
36 0.228097 e5m2 1 5
38 0.121183 e4m3 0 6
40 0.223447 e5m2 2 4
41 0.117631 e4m3 1 5
44 0.115716 e4m3 2 4
44 0.217276 e5m2 3 3
47 0.112486 e4m3 3 3
48 0.206713 e5m2 4 2
50 0.109289 e4m3 4 2
52 0.0884409 e5m2 5 1
53 0.0443644 e4m3 5 1
56 0.0161829 e4m3 6 0
56 0.0161829 e5m2 6 0
\end{filecontents*}

\begin{filecontents*}[overwrite]{references.bib}
@article{amsel2025polar,
  author  = {Amsel, Noah and Persson, David and Musco, Christopher and Gower, Robert M.},
  title   = {The {Polar} {Express}: Optimal Matrix Sign Methods and Their Application to the {Muon} Algorithm},
  journal = {arXiv preprint arXiv:2505.16932},
  year    = {2025}
}

@article{liu2025muon,
  author  = {Liu, Jingyuan and Su, Jianlin and Yao, Xingcheng and Jiang, Zhejun and Lai, Guokun and Du, Yulun and Qin, Yidao and Xu, Weixin and Lu, Enzhe and Yan, Junjie and others},
  title   = {Muon is Scalable for {LLM} Training},
  journal = {arXiv preprint arXiv:2502.16982},
  year    = {2025}
}

@article{du2026polarswitch,
  author  = {Du, Zhehang and Su, Weijie},
  title   = {Polar Switch: Fast {Muon} Orthogonalization with Polynomial Circuits},
  journal = {arXiv preprint arXiv:2610.26742},
  year    = {2026}
}

@article{huang2026orthogonalization,
  author  = {Huang, Hua},
  title   = {How Much Orthogonalization Does {Muon} Need?},
  journal = {arXiv preprint arXiv:2606.00371},
  year    = {2026}
}

@article{hariri2025quantize,
  author  = {Hariri, Mohsen and Luo, Alan and Chen, Weicong and Zhong, Shaochen and Zhang, Tianyi and Wang, Qifan and Hu, Xia and Han, Xiaotian and Chaudhary, Vipin},
  title   = {Quantize What Counts: More for Keys, Less for Values},
  journal = {arXiv preprint arXiv:2502.15075},
  year    = {2025}
}

@article{micikevicius2022fp8,
  author  = {Micikevicius, Paulius and Stosic, Dusan and Burgess, Neil and Cornea, Marius and Dubey, Pradeep and Grisenthwaite, Richard and Ha, Sangwon and Heinecke, Alexander and Judd, Patrick and Kamalu, John and Mellempudi, Naveen and Oberman, Stuart and Shoeybi, Mohammad and Siu, Michael and Wu, Hao},
  title   = {{FP8} Formats for Deep Learning},
  journal = {arXiv preprint arXiv:2209.05433},
  year    = {2022}
}

@article{micikevicius2018mixed,
  author  = {Micikevicius, Paulius and Narang, Sharan and Alben, Jonah and Diamos, Gregory and Elsen, Erich and Garcia, David and Ginsburg, Boris and Houston, Michael and Kuchaiev, Oleksii and Venkatesh, Ganesh and Wu, Hao},
  title   = {Mixed Precision Training},
  journal = {arXiv preprint arXiv:1710.03740},
  year    = {2018}
}

@article{dettmers2022optimizers,
  author  = {Dettmers, Tim and Lewis, Mike and Shleifer, Sam and Zettlemoyer, Luke},
  title   = {8-bit Optimizers via Block-wise Quantization},
  journal = {arXiv preprint arXiv:2110.02861},
  year    = {2022}
}

@inproceedings{gupta2015limited,
  author    = {Gupta, Suyog and Agrawal, Ankur and Gopalakrishnan, Kailash and Narayanan, Pritish},
  title     = {Deep Learning with Limited Numerical Precision},
  booktitle = {Proceedings of the 32nd International Conference on Machine Learning},
  year      = {2015},
  note      = {arXiv:1502.02551}
}

@inproceedings{gupta2018shampoo,
  author    = {Gupta, Vineet and Koren, Tomer and Singer, Yoram},
  title     = {Shampoo: Preconditioned Stochastic Tensor Optimization},
  booktitle = {Proceedings of the 35th International Conference on Machine Learning},
  year      = {2018},
  note      = {arXiv:1802.09568}
}

@article{higham1986polar,
  author  = {Higham, Nicholas J.},
  title   = {Computing the Polar Decomposition---with Applications},
  journal = {SIAM Journal on Scientific and Statistical Computing},
  volume  = {7},
  number  = {4},
  pages   = {1160--1174},
  year    = {1986}
}

@book{higham2008functions,
  author    = {Higham, Nicholas J.},
  title     = {Functions of Matrices: Theory and Computation},
  publisher = {SIAM},
  year      = {2008}
}

@book{golub2013matrix,
  author    = {Golub, Gene H. and Van Loan, Charles F.},
  title     = {Matrix Computations},
  edition   = {4},
  publisher = {Johns Hopkins University Press},
  year      = {2013}
}

@book{stewart1990perturbation,
  author    = {Stewart, G. W. and Sun, Ji-guang},
  title     = {Matrix Perturbation Theory},
  publisher = {Academic Press},
  year      = {1990}
}
\end{filecontents*}

\documentclass[10pt]{article}
\usepackage[T1]{fontenc}
\usepackage{times}
\usepackage[margin=1in]{geometry}
\usepackage{amsmath,amssymb,amsthm}
\usepackage{booktabs}
\usepackage{graphicx}
\usepackage{pgfplots}
\pgfplotsset{compat=1.18}
\usepgfplotslibrary{groupplots}
\usepackage{microtype}
\usepackage{array}
\setlength{\emergencystretch}{6em}
\usepackage[colorlinks=true,linkcolor=blue,citecolor=blue,urlcolor=blue]{hyperref}

\newtheorem{theorem}{Theorem}
\newtheorem{proposition}{Proposition}
\newtheorem{corollary}{Corollary}
\newtheorem{definition}{Definition}
\newtheorem{remark}{Remark}
\newtheorem{observation}{Observation}

\newcommand{\R}{\mathbb{R}}
\newcommand{\polar}{\operatorname{polar}}
\newcommand{\skewop}{\operatorname{skew}}
\newcommand{\tr}{\operatorname{tr}}
\newcommand{\norm}[1]{\left\lVert #1 \right\rVert}
\newcommand{\abs}[1]{\left\lvert #1 \right\rvert}
\newcommand{\e}[2]{\ensuremath{#1{\times}10^{#2}}}

\title{Precision-Budgeted Polynomial Iterations:\\
Sensitivity Certificates and Subspace Error Floors}
\author{Taehyeok Heo}
\date{}

\begin{document}
\maketitle

% --- BEGIN sections/00-abstract.tex ---
\begin{abstract}
Polynomial iterations that map a matrix to the polar factor (Newton--Schulz and minimax
variants, as used for Muon) run in bf16 and fp8, and the shipped designs carry
hand-tuned prescriptions with no account of what they cost at a given precision.  We
give that account in three separable pieces and spend it on the
three prescriptions that are wrong at four bits.
\textbf{Concretely, on the next Muon step:} raise the safety factor from $1.01$ to at
least $1.035$ --- it widens the tolerance of the certified window from $0.8\%$ to $4.8\%$
of its upper edge and lowers the hazard of one Muon call from $1$ in $1{,}600$ to $1$ in
$5{,}900$ on \texttt{digits} and from $1$ in $109$ to below $1$ in $23{,}000$ on
\texttt{diabetes} ($1{,}200$ steps, $25$ seeds each), which is margin and not immunity:
on \texttt{digits} it delays the escape, while keeping the state inside its window removes
it ($0$ in $210{,}000$ calls); stop at $K=4$--$5$, since the sixth step costs
eight bits for under $1.5\%$ of the window error and a four-step E4M3 update reaches
$1.3\times10^{-1}$ at 25 bits; and mix precision only inside a $\sigma_{\max}$-anchored
window with $\ell\le10^{-2}$, where alone it is not Pareto dominated, since at 40 bits a
four-step bf16 update errs by $2.4\times10^{-2}$ against $1.4\times10^{-1}$ for the
shipped bf16-cast, two-bf16-step hybrid.
The account: a computable sensitivity indicator $S_{\mathrm{pert}}(K)$ --- a prefix
ratchet of the coefficients' orbit sensitivity, needing no matrix and no iteration ---
forms with one calibration constant $c$ the law
$\mathrm{err} \approx E_K + c\,u\,S_{\mathrm{pert}}(K)$, $u = 2^{-p}$, whose median
predicted/measured ratio over 90 cells is $1.10$ (range $0.76$--$1.46$).  The floor is a subspace error, created at the first quantisation of a state whose spectrum still spans $[\ell,1]$; a first-order formula
predicts it within $10\%$ in 52 of 54 square cells and, in the two-term form a rectangular
update requires, within $0.6\%$ in all 500 momentum matrices of a Muon loop, so the floor
is $\Theta(u/\ell)$ and the earlier exponent $\ell^{-0.73}$ a chord.
Normalising the pair operator $[1/(\sigma_i+\sigma_j)]$ by $\sqrt d$ turns the floor's
family and dimension dependence into one number, $6.07$: a derived $2\sqrt2$ times a
profile factor ($1.31$ at the median over
180 cells, $0.98$--$2.10$) that the edge of a four-block variance-profile operator reproduces
to $1.2\%$ with no fitted constant, and that follows the decades a spectrum spans, not
$d$.
Inside real training (digits, diabetes; eleven seed-runs each) the loss is flat where the
update error is not: twelve bits and an update error of $0.73$ raise it by $1.65\times$
--- and what low precision breaks is survival, in a single
call of the iteration: one rounding lands the state past its window edge, the quintic's
gain then exceeds one ($2.2\times$ at $1.5$ times the edge, $9.3\times$ at twice it) and
the update runs to $10^{44}$, while clamping the state inside the window makes all fifty runs
finish and a $10\%$ learning-rate cut leaves seven of eleven failing.  In fp8 the floor matches the unlimited-exponent kernel's at the matched significand to
$1$--$3\%$, and online budget adaptation is a negative result: four spectrum-free sensors
cannot see the plateau ($69$ of $80$ matrices want $K=5$, the rest $K=4$).  No lower bound on accuracy versus
sensitivity holds.
\end{abstract}
% --- END sections/00-abstract.tex ---

% --- BEGIN sections/01-introduction.tex ---
\section{Introduction}

\paragraph{Motivation.}
Orthogonalizing the update matrix is standard in large-scale optimization: Muon and its
descendants replace the momentum matrix $M$ by its polar factor $UV^\top$
($M = U\Sigma V^\top$), and since an SVD per step is unaffordable the factor is computed
by a fixed polynomial iteration
\begin{equation}
X_{t} \;=\; p_t(X_{t-1}), \qquad
p_t(x) \;=\; a_t x + b_t x^3 + c_t x^5 , \qquad X_0 = M/\norm{M}_2 ,
\label{eq:iter}
\end{equation}
applied to the symmetric positive semidefinite form $XX^\top$.  The Polar Express
line of work~\cite{amsel2025polar} derives the coefficients minimizing
$\max_{\sigma\in[\ell,1]} |1 - p(\sigma)|$ for a given number of steps, where
$[\ell,1]$ is the spectrum of $X_0$ after normalization.  Optimality is a statement
about exact arithmetic; in $p$ significant bits neither part survives unqualified, and
the shipped designs compensate empirically with a safety factor $\approx1.01$, a cushion
$\approx0.02$ and a damped first step.

\paragraph{What to change in the next Muon step.}
Three prescriptions are wrong at four bits, each with its measured price.
\begin{enumerate}
\item \textbf{Safety $1.01 \to \ge 1.035$.}  At E5M2 and $K=6$ the shipped $1.01$
  loses 6 of 11 seed-runs to an overflow of the update within $400$ steps; $1.035$ loses
  none there and lowers the hazard of a single call $4$--$200\times$ over $1{,}200$ steps,
  for $0.4\%$ of the update's Frobenius norm.  That is margin, not immunity (\S\ref{sec:escape});
  the damping widens the excess a certified
  window tolerates from $0.8\%$ to $4.8\%$ of its upper edge, and one rounding beyond the
  edge lets the quintic's gain exceed one (\S\ref{sec:escape}).
\item \textbf{$K=6 \to 4$--$5$.}  The sixth step spends eight bits for under $1.5\%$ of
  the window error; $69$ of $80$ real momentum matrices want $K=5$ and the rest $K=4$;
  at 25 bits a four-step E4M3 update errs by $1.3\times10^{-1}$
  (\S\ref{sec:steps}).
\item \textbf{Mix precision only inside the window.}  The bf16-cast, two-bf16-step
  hybrid is Pareto dominated outside a $\sigma_{\max}$-anchored window with
  $\ell\le10^{-2}$: at 40 bits it errs by $1.4\times10^{-1}$ where a four-step bf16
  update errs by $2.4\times10^{-2}$, and is six times worse than bf16 on every real
  momentum matrix (\S\ref{sec:steps}).
\end{enumerate}

\paragraph{Contributions.}
\begin{enumerate}
\item \textbf{A computable sensitivity indicator and its calibration}
(Proposition~\ref{thm:cert}).  $S_{\mathrm{pert}}$ is a \emph{prefix ratchet} of the
orbit sensitivity of the coefficients: a scalar of $O(K)$ arithmetic, needing no
matrix.  The ratchet matters: the terminal sensitivity collapses to $1.0$ when the last
step is of Newton type and under-estimates the error by $66\times$.  With one constant
$c=0.142$ it forms $\mathrm{err} \approx E_K + c\,u\,S_{\mathrm{pert}}(K)$, whose median
ratio to measurement is $1.10$ over 90 cells.
\item \textbf{The floor is a subspace error, and it is invariant} (Theorems
\ref{thm:invariance}, \ref{thm:firstorder}).  The exact iteration preserves the polar
factor, so no step count, design or schedule reduces the floor; it is created at the first
quantisation of a state whose spectrum is still $[\ell,1]$.  A
first-order formula predicts it within $10\%$ in 52 of 54 square cells and, in the
two-term form rectangular matrices require, within $0.6\%$ in all 500 real momentum
matrices, and yields $\mathrm{err}=\Theta(u/\ell)$.
\item \textbf{The constant closes} (\S\ref{sec:profile}).  Normalising the pair operator
by $\sqrt d$ leaves $6.07$, a derived $2\sqrt2$ times a profile factor $g$ measured at
$1.31$ (median, 180 cells).  $g$ is the edge of a four-block
variance-profile operator built from the spectrum, with no fitted constant ($1.2\%$ of the
measurement at the median, $97\%$ of cells inside $\pm10\%$), and it follows the decades a
spectrum spans, not $d$.
\item \textbf{Two budget decisions} (\S\ref{sec:stopping}, \S\ref{sec:alloc}).  The
target rule $K(\varepsilon) = \min\{K : E_K + c\,u\,S_{\mathrm{pert}}(K) \le
\varepsilon\}$ holds on 62 cells (measured/target in $[0.02,1.15]$) and survives a
wrong dynamic range in one direction only (252 decisions); the graded schedule gains a
median $3.06\times$ where the first-slot premium exceeds two bits and is neutral on real
momentum ($0.9$ bits).
\item \textbf{Where the arithmetic ends: fp8} (\S\ref{sec:fp8}).  With the significand
matched, the fp8 floor equals the unlimited-exponent kernel's to $1$--$3\%$ over four
dynamic ranges, and the exponent window is a second, separable constraint with a computable
boundary; inside it, casting the start and running the first step in bf16 cures the
divergence, and outside it the same hybrid loses to dropping a step.
\item \textbf{What real training breaks, and how} (\S\ref{sec:realtrain},
\S\ref{sec:escape}).  The training loss is flat across four formats and update errors up
to $0.73$; what fails is survival, and it fails inside one call of the iteration rather
than in the training dynamics: the state's top leaves its certified window, the quintic
amplifies it ($2.2\times$ at $1.5$ times the edge, $9.3\times$ at twice), and the update
reaches $10^{44}$ in four steps.  Clamping the state into its window makes every long run
finish where a $10\%$ cut in the learning rate leaves 7 of 11 failing, and the curvature loop
gain stays below one ($\le0.83$ in five of six trajectories), which excludes an
edge-of-stability story.
\item \textbf{Three negative results.}  Accuracy does not require sensitivity (the error
falls $500\times$ while the sensitivity is non-monotone); the law does not design (at
bf16 the shipped point is within $1.11\times$ of the best of 35 designs and the law's own
argmin is $1.58\times$ \emph{worse}); and no spectrum-free sensor predicts which matrices
gain from precision (Spearman $\le0.33$, \S\ref{sec:online}).
\end{enumerate}
% --- END sections/01-introduction.tex ---

% --- BEGIN sections/02-preliminaries.tex ---
\section{Preliminaries}
\label{sec:prelim}

\paragraph{Setting.}
Let $M \in \R^{m\times n}$ have full column rank, $m\ge n$, and let
$X_0 = M/\norm{M}_2$ with $\sigma(X_0) \subset [\ell,1]$; the target is the polar factor
$\polar(M) = UV^\top$, $M = U\Sigma V^\top$.  A Muon update has this shape with $m = H$,
$n = D_{\mathrm{in}}$; $d=m=n$ in the square experiments and \S\ref{sec:realmomentum}
reports the rectangular case.  A step with coefficients $(a,b,c)$ is
\begin{equation}
X \mapsto aX + b\,(XX^\top)X + c\,(XX^\top)^2 X ,
\label{eq:step}
\end{equation}
$P_K$ is the composition of $K$ steps, a design is the sequence $(a_t,b_t,c_t)_{t\le K}$,
and the shipped family is the minimax quintic sequence of~\cite{amsel2025polar} with
cushion $0.02$ and safety $1.01$ unless stated.
$E_K = \max_{\sigma\in[\ell,1]}|p_K(\cdots p_1(\sigma)) - 1|$ is the exact-arithmetic
error, a scalar because the exact iteration can be simulated on $[\ell,1]$ alone.

\paragraph{Kernel model.}
We model storage, not the arithmetic: $x$ is stored as
$q(x,p) = \mathrm{ldexp}(\mathrm{round}(\mathrm{mant}(x)2^{p-1})2^{1-p},\allowbreak
\mathrm{exp}(x))$, i.e.\ $p$ significant binary digits, with relative rounding error at
most $2^{1-p}$ and rms $u/\sqrt3$ for $u = 2^{-p}$, and one step is
\begin{equation}
A = XX^\top, \qquad B = q\bigl(bA + cA^2,\, p\bigr), \qquad
X \leftarrow q\bigl(aX + BX,\, p\bigr),
\label{eq:kernel}
\end{equation}
with products accumulated in float64, a wide accumulator: the standard picture of a
low-precision kernel, reproducing what matters (a floor $\propto u$, saturation in $K$,
divergence) with every claim attributable to one equation.  A schedule $(p_0,\dots,p_K)$
gives the precision of the input cast and of each step.

\paragraph{Orbit sensitivity and its ratchet.}
For a design $(p_t)_{t\le K}$ and $\sigma\in[\ell,1]$ let $x_0=\sigma$ and\begin{equation}
S_K(\sigma) = \max_{1\le j \le K}\ \prod_{l=j}^{K} \bigl|p_l'(x_{l-1})\bigr|,
\qquad S(K) = \max_{\sigma\in[\ell,1]} S_K(\sigma).
\label{eq:S}
\end{equation}
$S_K(\sigma)$ is the largest amplification a rounding injected at step $j$ suffers on its
way to the output; the product form is exact because each $p_l$ is differentiable and the
orbit is one-dimensional.  $S(K)$ is a scalar of the coefficients: on a 400-point grid it
costs $O(K)$ arithmetic per point, needs no matrix and no SVD.

\begin{definition}[Prefix ratchet sensitivity indicator]
\label{def:ratchet}
$S_{\mathrm{pert}}(K) = \max_{1\le t\le K} S(t)$.
\end{definition}

The ratchet is what makes $S$ usable.  A Newton-type final step has $p'(1)=0$, so the
terminal factor of $S_K$ vanishes and $S(K)$ collapses to $1.0$ exactly where the floor
dominates: at $\ell=10^{-3}$, $K=7$, $p=8$ it predicts $3.9\times10^{-3}$ against a measured
$2.58\times10^{-1}$, a $66\times$ under-estimate, while the off-diagonal gain of that step
is $1/2$ --- which the prefix maximum retains (Fig.~\ref{fig:ratchet}).

\paragraph{Decomposition.}
We separate $\mathrm{err} = \norm{X_K - \polar(M)}_2$ into the singular-value part
$\sigma_{\mathrm{err}} = \norm{X_K - \polar(X_K)}_2$ and the subspace part
$\mathrm{sub}_{\mathrm{err}} = \norm{\polar(X_K) - \polar(M)}_2$; over 12 cells the
subspace part is $86$--$100\%$ of the total (\S\ref{sec:invariance} explains why).
% --- END sections/02-preliminaries.tex ---

% --- BEGIN sections/03-certificate-subspace.tex ---
\section{The Certificate and the Subspace Floor}
\label{sec:cert}

\subsection{The sensitivity indicator and the two-term law}
\label{sec:thmcert}

\begin{proposition}[Two-term law: an empirical predictive model]
\label{thm:cert}
Let the design have exact error $E_K$ and prefix ratchet $S_{\mathrm{pert}}(K)$ as
in \S\ref{sec:prelim} and Definition~\ref{def:ratchet}, and let the state and the
intermediate $B$ be stored with $p$ significant bits, $u = 2^{-p}$, in a kernel whose
exponent range does not bind.  Then
\begin{equation}
\mathrm{err} \;\approx\; E_K \;+\; c\,u\,S_{\mathrm{pert}}(K) \;+\; O(u^2),
\qquad c = 0.142,
\label{eq:cert}
\end{equation}
to within a factor $1.6$ in all 192 conditions of the design family (worst deviation
$1.60$, CI $[1.59,1.60]$), and the error after the stopping point is within
$0.76$--$1.46\times$ of the prediction over a 90-cell sweep of dynamic ranges
$\ell\in[10^{-4},10^{-1}]$, precisions $p\in[4,24]$ and $d\in\{64,128\}$ (median ratio
$1.10$, CI $[1.07,1.14]$).
\end{proposition}

We call this a proposition deliberately: $c$ is a calibration and the expression a law
fitted and tested, not a theorem with a controlled remainder.  A
theorem of the same shape exists --- $(2K+1)\,u\,S_{\mathrm{pert}}$ bounds the
singular-value channel with no calibration, and the subspace channel with the
off-diagonal gain in place of $p'$ (Proposition~\ref{prop:ratchet}) --- and the gap
between them is the calibration: the off-diagonal gain of a step is not dominated by its
derivative, so $c=0.142$ compresses a bound by two orders of magnitude.  The
first term is the exact error a design optimizes, the second the floor; the second is flat
in $K$ and the first falls geometrically, so their sum has an interior minimum, where the
shipped step counts sit.

\begin{table}[t]
\centering
\caption{Certificate against measurement.  Shipped design, $\ell=10^{-3}$,
$d=64$, $K=6$, $E_6 = 1.623\times10^{-3}$, $S_{\mathrm{pert}}(6) = 609.6$; $q$ is
the input cast and both intermediate roundings.  Ratio $=$ predicted/measured.}
\label{tab:cert}
\begin{tabular}{rrrr}
\toprule
$p$ & measured & predicted & ratio \\
\midrule
6  & 1.66          & 1.35          & 0.82 \\
7  & \e{6.15}{-1}  & \e{6.78}{-1}  & 1.10 \\
8  & \e{2.63}{-1}  & \e{3.40}{-1}  & 1.29 \\
9  & \e{1.98}{-1}  & \e{1.71}{-1}  & 0.86 \\
10 & \e{8.36}{-2}  & \e{8.62}{-2}  & 1.03 \\
11 & \e{3.39}{-2}  & \e{4.39}{-2}  & 1.29 \\
12 & \e{1.89}{-2}  & \e{2.28}{-2}  & 1.20 \\
14 & \e{5.14}{-3}  & \e{6.91}{-3}  & 1.34 \\
16 & \e{2.31}{-3}  & \e{2.94}{-3}  & 1.28 \\
20 & \e{1.64}{-3}  & \e{1.71}{-3}  & 1.04 \\
24 & \e{1.62}{-3}  & \e{1.63}{-3}  & 1.00 \\
\bottomrule
\end{tabular}
\end{table}

Three regimes fit (Table~\ref{tab:cert}, ratios $0.82$--$1.34$): the floor
dominates at $p\le8$, the exact error at $p\ge20$; at $p=4$ the model predicts
non-convergence and the measurement diverges.

\subsection{The exact iteration preserves the singular subspace}
\label{sec:invariance}

\begin{theorem}[Invariance]
\label{thm:invariance}
Let $p(x) = ax+bx^3+cx^5$ with real coefficients and
$P(X) = aX + b(XX^\top)X + c(XX^\top)^2X$.  If $X = U\Sigma V^\top$ is full rank
and $p(\sigma_i)\neq0$ for all $i$, then $P(X) = U\,p(\Sigma)\,V^\top$ and hence
$\polar(P(X)) = UV^\top = \polar(X)$.  The same holds for any composition $P_K$ of
such steps with different coefficients.
\end{theorem}

\begin{proof}
$XX^\top = U\Sigma^2U^\top$, so $(XX^\top)X = U\Sigma^3V^\top$ and
$(XX^\top)^2X = U\Sigma^5V^\top$; summing gives
$U(a\Sigma+b\Sigma^3+c\Sigma^5)V^\top = Up(\Sigma)V^\top$, whose polar factor is
$UV^\top$ whenever all $p(\sigma_i)$ are positive.  Induct over the composition.
\qed
\end{proof}

\begin{corollary}[Why the floor is flat]
\label{cor:flat}
For the $p$-bit orbit $X_t = q(P_t(X_{t-1}),p)$, $\mathrm{sub}_{\mathrm{err}} =
\norm{\polar(X_K) - \polar(X_0)}$: the exact part of the iteration contributes
exactly zero to the subspace error.  Hence (i) the subspace floor does not fall with
$K$, (ii) it does not depend on the design, and (iii) it does not depend on how bits
are spread over the steps --- it is decided by the first quantisation of a state
whose spectrum spans $[\ell,1]$.
\end{corollary}

In exact arithmetic the polar factor stays fixed to $4\times10^{-14}$ at $\ell=10^{-3}$
over the whole orbit, and a known rotation between the two trailing singular directions is
transported unchanged through six steps: the earlier contraction--reinjection picture
bounds the singular-value part only.

\subsection{The size of the floor: a first-order formula}
\label{sec:firstorder}

By Corollary~\ref{cor:flat} the subspace error is the response of the polar factor
to the perturbation $E = q(\hat M,p) - \hat M$ of the start.

\begin{theorem}[Polar response to a perturbation]
\label{thm:firstorder}
Let $\hat M = U\Sigma V^\top \in \R^{m\times n}$ have full column rank, $m\ge n$,
$\sigma_{\min}(\hat M) = \ell$, let $E$ be a perturbation and $F = U^\top E V$.
Then
\begin{equation}
\begin{aligned}
\polar(\hat M + E) - \polar(\hat M)
=\;& \underbrace{U\left[\frac{F_{ij}-F_{ji}}{\sigma_i+\sigma_j}\right]_{ij} V^\top}_{\text{rotation inside } \operatorname{range}(U)}
\;+\; \underbrace{(I_m - UU^\top)\,E\,(\hat M^\top\hat M)^{-1/2}}_{\text{rotation of } \operatorname{range}(U)}
\;+\; O(\norm{E}^2).
\end{aligned}
\label{eq:firstorder}
\end{equation}
For square $\hat M$ the second term vanishes identically and
\eqref{eq:firstorder} reduces to the single term used in the synthetic tests; for
$m<n$ the same statement holds with $U$ and $V$ exchanged.
\end{theorem}

\begin{proof}
Perturb $\Sigma \to \Sigma + tF$ and write the first-order SVD as
$(I+t\Omega_u)(\Sigma+tD)(I+t\Omega_v)^\top$ with $\Omega_u,\Omega_v$ skew.  Matching
off-diagonal entries gives $(\Omega_u-\Omega_v)_{ij} = (F_{ij}-F_{ji})/(\sigma_i+\sigma_j)$,
and the polar factor is $(I+t\Omega_u)(I+t\Omega_v)^\top$; this covers the part of $E$
inside $\operatorname{range}(U)\times\R^n$.  For the complement, the differential of
$\polar(A) = A(A^\top A)^{-1/2}$ has a summand inside $\operatorname{range}(U)$ and a
first that contributes $(I_m-UU^\top)E(\hat M^\top\hat M)^{-1/2}$.  \qed
\end{proof}

Three consequences matter.  The \emph{symmetric part of a perturbation does not move the
polar factor}, and the in-subspace gain peaks at $1/(2\ell)$, so quantisation is expensive
only while the spectrum is wide.  The \emph{projected term is not damped pair by pair}:
it carries $\norm{E}/\ell$ in any direction and is absent from every square test.  And the first term is the fixed operator $T(e_{ij})=(\sigma_i+\sigma_j)^{-1}(e_{ij}-e_{ji})$,
answered by $2\sqrt2\,\rho\norm{W}_F/\sqrt d$ for an exchangeable perturbation
(Appendix~\ref{app:constant}).

Across the grid the prediction is within $10\%$ in 52 of 54 cells
(Table~\ref{tab:firstorder}) and $\norm{E}_2/u$ is measured $0.46$--$0.89$.  The grid is square and a Muon update is not: on the 500 real momentum
matrices the projected term is larger than the in-subspace one by a median $2.77$
(\S\ref{sec:loop}), so the floor is set by the rotation of the column space, not by the
damping of the pair weights.

\begin{corollary}[Scale, and what the exponent $0.73$ was]
\label{cor:scale}
$\mathrm{sub}_{\mathrm{err}} \le (\norm{\skewop F}_2 + \norm{E}_2)/\ell$, so
$\mathrm{sub}_{\mathrm{err}} = \Theta(u/\ell)$ while the response is unsaturated and
$\norm{E}_2 = \Theta(u)$: both terms of \eqref{eq:firstorder} are bounded by
$\norm{E}_2/\ell$, the first because $1/(\sigma_i+\sigma_j)\le1/(2\ell)$, the second
because $\norm{(\hat M^\top\hat M)^{-1/2}}_2 = 1/\ell$.
\end{corollary}

The measured $C(\ell)$ follows \eqref{eq:firstorder} over four decades (predicted over
measured $0.49$--$1.06$ at $d=64$, $p=12$), but its log--log fit is not a power law: the
local slope rises from $0.2$ to $1.17$, and $-0.73$ is the chord of that crossover.
Appendix~\ref{app:bound} has the matrix-free bound and the impossibility of a spectral one.

\subsection{Consequences for the budget and how they fail}
\label{sec:consequences}

The floor is the sum of the responses to the per-slot roundings, each amplified at the
spectrum the state has when it happens, so allocation reduces to the per-slot
gains $\kappa_t$ (\S\ref{sec:alloc}); a step pays only while the exact error falls faster
than a floor that does not fall with $K$ --- and on real spectra one step before the
shipped count (\S\ref{sec:steps}).

\label{sec:failures}
Four failures are kept in the text: saturation (at $\mathrm{sub}_{\mathrm{err}}=2$,
\eqref{eq:firstorder} over-predicts by $2.9\times$), a calibrated constant with a
$2.1$--$14\times$ loose matrix-free bound (Appendix~\ref{app:bound}), no design lever
(Appendix~\ref{sec:design}), and a loss that is not a meter (\S\ref{sec:realtrain}).
% --- END sections/03-certificate-subspace.tex ---

% --- BEGIN sections/04-stopping.tex ---
\section{Stopping: from a Crossing to a Target}
\label{sec:stopping}

The law \eqref{eq:cert} estimates the error at a given $K$; a practitioner needs a
decision, and there are two ways to read it as a rule.  The first asks where the precision
term overtakes the exact term,
\begin{equation}
K_{\mathrm{cross}}(u) \;=\; \min\{K:\ E_K \le u\,S_{\mathrm{pert}}(K)\},
\label{eq:cross}
\end{equation}
the reading the two-term structure invites; as a stopping time it is also wrong by up to
three orders of magnitude in error (Table~\ref{tab:cross}, appendix).  The second asks which $K$ first meets a stated accuracy,
\begin{equation}
K(\varepsilon) \;=\; \min\{K:\ J(K) \le \varepsilon\},
\qquad J(K) = E_K + c\,u\,S_{\mathrm{pert}}(K),
\label{eq:target}
\end{equation}
and it is the one we recommend and verify.

\begin{theorem}[Stopping at a target accuracy]
\label{thm:stop}
Let $J(K) = E_K + c\,u\,S_{\mathrm{pert}}(K)$ with the calibrated $c$ and $K_{\max}$ the
largest step count the caller will run.  If $\varepsilon > \min_{K \le K_{\max}} J(K)$
then $K(\varepsilon)$ exists and is finite; if $\varepsilon$ is below the predicted floor
the rule returns nothing, and no step count exists.  At the returned step count
$\mathrm{err}(K(\varepsilon)) \le 1.16\,\varepsilon$ in all 62 cells tested (median
$0.45\,\varepsilon$), and it is below $\varepsilon$ in 60 of 62.
\end{theorem}

Existence is immediate; the content is the last sentence, measured over
$d\in\{64,128\}$, three $\ell$, four $p$ and eight
targets from $0.3$ to $10^{-4}$ (Table~\ref{tab:stop}); at $\ell=10^{-3}$, $p=10$ the
predicted floor, $0.0845$, takes the refusal branch.

\begin{table}[t]
\centering
\caption{Target-accuracy stopping rule \eqref{eq:target} against measurement.  Eight
representative rows of the 62 defined cells.  Ratio $=$ measured error / target; the
refusal branch is taken exactly when the target is below the predicted floor
$\min_K J(K)$.}
\label{tab:stop}
\begin{tabular}{rrrr rc}
\toprule
$\ell$ & $p$ & $\varepsilon$ & $K(\varepsilon)$ & measured & ratio \\
\midrule
$10^{-2}$ & 12 & $0.3$  & 3 & \e{2.80}{-1} & 0.93 \\
$10^{-2}$ & 12 & $0.1$  & 4 & \e{1.67}{-2} & 0.17 \\
$10^{-2}$ & 12 & $0.01$ & 5 & \e{3.46}{-3} & 0.35 \\
$10^{-3}$ & 12 & $0.3$  & 5 & \e{1.36}{-1} & 0.45 \\
$10^{-3}$ & 12 & $0.1$  & 6 & \e{1.89}{-2} & 0.19 \\
$10^{-3}$ & 16 & $0.003$& 6 & \e{2.21}{-3} & 0.74 \\
$10^{-4}$ & 16 & $0.3$  & 7 & \e{4.50}{-2} & 0.15 \\
$10^{-4}$ & 16 & $0.03$ & 8 & \e{9.55}{-3} & 0.32 \\
\midrule
\multicolumn{3}{r}{summary, 62 cells} & \multicolumn{2}{c}{ratio $= 0.020$--$1.153$, median $0.451$} \\
\multicolumn{3}{r}{within $1.5\times$} & \multicolumn{2}{c}{62/62} \\
\multicolumn{3}{r}{below the target} & \multicolumn{2}{c}{60/62} \\
\bottomrule
\end{tabular}
\end{table}

The rule is not monotone in $\varepsilon$: over the 44 consecutive
target pairs $K(\varepsilon)$ moves in the expected direction 26 times, because $J$ is not
monotone in $K$; both members of an inverted pair still meet their own targets, and a
$3.5\times$ over-estimate of $c$ changes no selected $K$.

\paragraph{Why the crossing rule is the wrong stopping time.}
Equation~\eqref{eq:cross} does describe where the floor takes over, at the $u$-scale
\eqref{eq:cert} predicts, but it does not locate the argmin: on the frozen real momentum
spectra ($\ell=10^{-3}$, $K$ swept to 10) it stops one to three steps early --- at $p=8$
where the error is $9.1\times$ its value at $K=6$, and $926\times$ in the wide window at
$p=24$ --- because it asks where the two terms are equal in size, where the sum is still
dominated by the falling term, and the error is minimized where they are equal in
\emph{slope}; stopping late is cheap ($6.237\times10^{-2}$ at $K=10$ against
$6.229\times10^{-2}$ at $K=6$).

\paragraph{A lower bound on the dynamic range.}
\label{sec:ellsens}
The last input is $\ell$, which a caller does not have.  With the design and the rule
given $\hat\ell=\alpha\ell$, over 36 cells ($d\in\{64,128\}$,
three $\ell$, $p\in\{8,12\}$) at three targets each
($\varepsilon=r\min_K J(K)$, $r\in\{1.2,3,10\}$; Table~\ref{tab:ell}),
under-estimating is safe and costs arithmetic: at $\alpha\le0.5$ no target is missed, the
error stays below $0.99\,\varepsilon$, and the rule either returns a step count within
three steps of the right one or refuses a target it could have met (8 to 16 of 36 cells,
all at $r=1.2$).  Over-estimating is neither safe nor visible: at $\alpha=10$, 24 targets
are missed with errors up to $244\,\varepsilon$ and the caller's own prediction
under-states the measurement by up to $1275\times$ ($\e{6.1}{-4}$ against $0.782$ at
$\ell=10^{-2}$, $p=12$).  A design fitted to $[\hat\ell,1]$ leaves the polynomial
outside its minimax interval below $\hat\ell$, so the orbit converges to the wrong point.
So $\sigma_{\min}$ should be passed as a lower bound --- a power iteration count with a
safety factor.
% --- END sections/04-stopping.tex ---

% --- BEGIN sections/05-bit-allocation.tex ---
\section{Allocation: the Floor Is a Maximum, Not a Sum}
\label{sec:alloc}

Given $K$ steps and $B$ bits for the state and the intermediate, which slot deserves
them?  Each rounding contributes an independent response, and the
responses combine as a maximum rather than a sum.

\begin{proposition}[Max form]
\label{prop:max}
Let $\kappa_t = \mathrm{err}(\text{only slot } t \text{ at } p)/u$ be the response
to the rounding of slot $t$ alone.  Then
\begin{equation}
\mathrm{err} \;\approx\; \max_{0\le t\le K} \kappa_t\,2^{-p_t},
\qquad\text{so}\quad
\min_{\sum_t p_t \le B}\ \mathrm{err}
\ \text{ is solved by }\
p_t^\star = \log_2 \kappa_t + \tfrac{1}{n}\bigl(B - \textstyle\sum_s \log_2\kappa_s\bigr),
\label{eq:maxform}
\end{equation}
with $n = K+1$.  The gain over uniform precision at equal total bits is at most
$\max_t \kappa_t / (\prod_t \kappa_t)^{1/n}$, and exactly one when the $\kappa_t$ are
equal.
\end{proposition}

The support for \eqref{eq:maxform} is the pair of ratios
$\mathrm{err}/\max_t(\text{channel}) = 1.00$--$1.21$ and
$\mathrm{err}/\sum_t(\text{channel}) = 0.71$--$1.00$ on synthetic cells, $1.42$--$1.50$
against $0.11$ on the real rectangular ones.

\paragraph{Across the steps, and across the two operands.}
By Corollary~\ref{cor:flat} the per-slot gains follow the conditioning along the orbit.
Table~\ref{tab:profile} (appendix) shows the profile falling by $\approx2.2$ per slot to a
$1.1u$ plateau, with a first-slot premium of $3.35$ bits at $\ell=10^{-2}$ and $5.77$ at
$10^{-3}$; across the state and the intermediate the bottleneck is whichever operand is
coarser, $236\times$ for a $16$-bit imbalance (Table~\ref{tab:k1},
Appendix~\ref{app:refute}).  The mechanism is the spectrum, not the step index: on the
real momentum matrices the per-slot responses at $p=8$, $K=6$ are $2.36u$, $1.57u$,
$1.37u$ and a $1.27u$ plateau, a premium of $0.90$ bits rather than $5.8$.\begin{table}[t]
\centering
\caption{Equal total bits, uniform versus the graded schedule
\eqref{eq:maxform} ($d=64$), with the closed form of Table~\ref{tab:profile}.
Entries are $(p_0,p_1,\dots,p_K)$; gain is uniform error over graded error.}
\label{tab:schedule}
\begin{tabular}{rrrrccc}
\toprule
$\ell$ & $K$ & $B$ & uniform $p$ & schedule $(p_0,\dots,p_K)$ & error & gain \\
\midrule
$10^{-2}$ & 6 & 56 & 8 & $(10,9,9,7,7,7,7)$ & \e{2.54}{-2} & $2.02$ \\
$10^{-3}$ & 6 & 70 & 10 & $(14,12,11,10,8,8,7)$ & \e{2.03}{-2} & $4.12$ \\
$10^{-3}$ & 8 & 90 & 10 & $(14,13,11,10,9,8,8,8,8)$ & \e{1.32}{-2} & $6.29$ \\
$10^{-4}$ & 8 & 90 & 10 & $(16,14,12,11,9,8,7,7,6)$ & \e{4.30}{-2} & $15.8$ \\
$10^{-4}$ & 6 & 84 & 12 & $(17,15,13,12,10,9,8)$ & \e{3.99}{-1} & $1.22$ \\
\bottomrule
\end{tabular}
\end{table}

The solution is inverse waterfilling \eqref{eq:maxform}.  Fitting Table~\ref{tab:profile} as $\kappa_t \approx \max\{(\kappa/\ell)\rho^t, 1.1\}$
with $\kappa\approx0.06$, $\rho\approx0.45$ gives a rule with no matrix: the first slot
takes $\log_2(1/\ell)-4.2$ bits more than the tail and each later slot $1.2$ less.  Over
48 cells ($d\in\{64,128\}$,
$\ell\in\{10^{-2},10^{-3},10^{-4}\}$, $K\in\{6,8\}$, four budgets) the schedule gains a
median $3.06\times$ over uniform precision (CI $[2.24,4.76]$, range
$1.02$--$18.1\times$), and its largest-remainder rounding is exactly optimal in 44 of the
48 cells.  The last row of Table~\ref{tab:schedule} is the exception ($1.22$): allocation
is a remedy for the precision floor only.

\paragraph{Where the grading pays.}
\label{sec:premium}
Those 48 cells are synthetic logspace spectra; on 60 of the 500 real momentum matrices
($K\in\{4,6,8\}$, four budgets) the graded schedule buys nothing --- median gain $1.000$,
$p_{10}$--$p_{90}$ $0.887$--$1.309$.  The reason is granularity: one bit from a tail slot to the first pays only if the premium
clears two bits, and at the real range it is $0.9$.  A sweep of the same $128\times32$ shape (Table~\ref{tab:premium}, appendix)
gives $1.05$--$1.19\times$ at a premium of $1.6$ bits, $3.2$--$4.9\times$ at $4.4$ and
$6.4$--$15\times$ at $7.4$: allocation is a remedy for dynamic range and nothing at the
range Muon produces on a small hidden layer, which needs one fewer step instead
(\S\ref{sec:steps}).

\paragraph{Online adaptation does not see the plateau (negative).}
\label{sec:online}
If the per-slot gains are unknown, the natural move is to read them from the iteration.
Four sensors that cost no extra matrix product --- the progress of the state, the noise the
last step injected, the defect
$\norm{X_tX_t^\top-I}_F/\norm{X_tX_t^\top}_F$, and a two-iteration $\lambda_{\max}$ with a
Hutchinson trace --- plus an exit rule that widens a slot while the progress exceeds its own
noise change nothing: all three settings keep the bits of the uniform format they start
from and land on its error.  The reason is in the error columns of
\S\ref{sec:steps}: on the full spectrum every design errs by $1.00$--$1.65$
at every precision from $24$ to $168$ bits, because the median dynamic range is
$2\cdot10^{-10}$ and the tail is unreachable in six steps, so a global statistic is
dominated by modes the design never moves.

Nor is there anything for a per-matrix rule to allocate: every one of the 80 matrices is
within $1.15\times$ of its best error at $K=4$ or $5$, the best Spearman of a measurable
feature with the E4M3-over-E5M2 gain is $0.33$, and giving the wider format to the best
quarter by any feature recovers $26$--$27\%$ of the gain, as a random quarter does: the
schedule is a closed form for synthetic windows, not an online algorithm.
% --- END sections/05-bit-allocation.tex ---

% --- BEGIN sections/06-experiments.tex ---
\section{Experiments}
\label{sec:exp}

Every number here and in \S\ref{sec:cert}--\S\ref{sec:alloc} is a field of a JSON file
written by a script of Appendix~\ref{app:repro}; every median the abstract repeats carries
a bootstrap interval (Table~\ref{tab:ci}).  The synthetic family is $M = U\Sigma V^\top$
with random orthogonal $U,V$ and $\Sigma$ logspace from $1$ to $\ell$: $d\in\{32,64,128\}$,
$\ell\in[3\cdot10^{-5},3\cdot10^{-1}]$, $p\in\{4,\dots,24\}$, the shipped quintic unless
stated.  Two sets leave it: the momentum of a real Muon loop (500
matrices, \S\ref{sec:realmomentum}) and training itself (\S\ref{sec:realtrain}) ---
five-layer and residual MLPs on \texttt{digits}, a four-layer MLP on \texttt{diabetes},
Frobenius input normalisation, per-sample LayerNorm, momentum $0.95$ Nesterov,
$\ell_0=10^{-2}$, $K=6$ --- where the error is read on the design window, the directions
with $\sigma\ge0.05\norm{M}_F$, because the full-spectrum error is truncation.

\subsection{The floor, the certificate, and their scaling}

The floor is flat in $K$ and linear in $u$: between $p=10$ and $p=12$ the
$\ell=10^{-2}$ floor falls $3.75\times$ for a $4\times$ fall in $u$; and divergence is not
only a precision effect --- the undamped minimax sequence on a $d=64$ matrix reaches
$1.9\times10^{6}$ at bf16, $K=7$, where fp16 gives $0.0124$.

The certificate is only as good as its indicator, and the indicator's terminal value is
not usable (Fig.~\ref{fig:ratchet}).  At $\ell=10^{-3}$ the last step of the $K=7$ design
has $p'(1)=0$: $S(7)=2.5\times10^{-3}$ predicts $5.5\times10^{-4}$ where $0.258$ is
measured, while the ratchet stays at $610$ and puts the law at $0.34$; the floor moves
$1.5\times$ from $K=5$ to $K=7$ while the terminal sensitivity falls $609$-fold.

\begin{figure}[t]
\centering
% --- BEGIN figures/fig1-ratchet-and-floor.tex ---
% The ratchet: the sensitivity of every prefix, and what it predicts.
% Coordinates: figures/data/fig1-*.dat, written by data/make_figure_data.py from
% pilot24_results.json (c_design, L0.001_K*) and pilot5c_results.json (grid_K*, shipped).
\begin{tikzpicture}
\begin{groupplot}[
  group style={group size=2 by 1, horizontal sep=1.6cm},
  width=0.47\linewidth, height=4.3cm,
  tick label style={font=\footnotesize},
  label style={font=\footnotesize},
  title style={font=\footnotesize},
  legend style={font=\scriptsize, draw=none, fill=none},
  axis line style={thick}, grid=major, grid style={densely dotted, black!18},
]
\nextgroupplot[
  ymode=log, xmode=linear,
  xlabel={prefix $t$}, ylabel={$S(t)$},
  ymin=1e-3, ymax=3e3, xtick={1,2,3,4,5,6,7},
  title={(a) prefixes of the $K=7$ design, $\ell=10^{-3}$},
]
\addplot[black, thick, mark=*, mark size=1.7pt]
  table[x=t, y=S] {figures/data/fig1-ratchet-prefix.dat};
\addplot[gray, dashed, thick] coordinates {(5,609.6) (5,0.0025)};
\node[font=\scriptsize, anchor=west] at (axis cs:5.15,120) {terminal step};
\node[font=\scriptsize, anchor=west] at (axis cs:5.15,2.5) {$p_K'(1)=0$};
\node[font=\scriptsize, anchor=south west] at (axis cs:2.1,900) {ratchet $=610$};

\nextgroupplot[
  ymode=log,
  xlabel={step budget $K$}, ylabel={error},
  xmin=4.6, xmax=7.4, xtick={5,6,7}, ymin=2e-4, ymax=1,
  title={(b) the floor, and the two predictions},
  legend pos=south west, legend cell align=left,
]
\addplot[black, thick, mark=square*, mark size=1.7pt]
  table[x=K, y=meas8] {figures/data/fig1-ratchet-floor.dat};
\addplot[black, thick, mark=o, mark size=1.9pt, only marks]
  table[x=K, y=meas11] {figures/data/fig1-ratchet-floor.dat};
\addplot[gray, thick, dashed]
  table[x=K, y=pred8ratchet] {figures/data/fig1-ratchet-floor.dat};
\addplot[gray, thick, densely dotted, mark=triangle*, mark size=1.7pt]
  table[x=K, y=pred8term] {figures/data/fig1-ratchet-floor.dat};
\legend{measured $p=8$, measured $p=11$,
  $0.142\,u\,S_{\mathrm{pert}}$, $0.142\,u\,S(K)$}
\end{groupplot}
\end{tikzpicture}
% --- END figures/fig1-ratchet-and-floor.tex ---

\caption{The ratchet against the terminal sensitivity (\texttt{pilot24}, \texttt{pilot5c}).
(a) The sensitivity of every prefix of the shipped $K=7$ design at $\ell=10^{-3}$: it peaks
at $610$ at $t=5$ and the sixth step returns it to $2.5\times10^{-3}$.  (b) The measured
floor at $d=64$ (markers) against $0.142\,u\,S_{\mathrm{pert}}(K)$, which is flat, and
$0.142\,u\,S(K)$, which collapses; $p=11$ is $2.8$--$7.8\times$ lower with the same shape.}
\label{fig:ratchet}
\end{figure}

\subsection{Where the floor lives: the real momentum spectra}
\label{sec:realmomentum}

A two-layer MLP (input 32, hidden 128, 10 classes, synthetic teacher labels) trains for
500 steps with Muon on the hidden weight; we record all 500
momentum matrices ($128\times32$) in exact arithmetic and replay each through the $p$-bit
kernel (median dynamic range $0.150$).

\begin{table}[t]
\centering
\caption{Theorem~\ref{thm:firstorder} on the 500 momentum matrices of a Muon loop
($128\times32$, $K=8$, design window $0.150$).  ``in-subspace'' uses only the first
term of \eqref{eq:firstorder}; ``two-term'' adds
$(I-UU^\top)E(\hat M^\top\hat M)^{-1/2}$.  Direction error is
$\norm{\Delta_{\mathrm{pred}}-\Delta_{\mathrm{true}}}_2/\norm{\Delta_{\mathrm{true}}}_2$.}
\label{tab:real500}
\begin{tabular}{rcccccc}
\toprule
$p$ & in-subspace & two-term & within $10\%$ & direction error & $\norm{\text{term}_2}/\norm{\text{term}_1}$ & median $C\ell$ \\
\midrule
8  & $0.352$ & $1.000\ [0.994,1.006]$ & $500/500$ & $0.011$ & $2.77$ & $0.351$ \\
11 & $0.352$ & $1.000\ [0.999,1.001]$ & $500/500$ & $0.001$ & $2.77$ & $0.350$ \\
16 & $0.353$ & $1.000\ [1.000,1.000]$ & $500/500$ & $0.000$ & $2.75$ & $0.348$ \\
\bottomrule
\end{tabular}
\end{table}

With both terms the prediction is within $0.6\%$ in all 500 cells at $p=8$ and $0.1\%$ at
$p=16$; with the in-subspace term alone it is $0.35\times$ the measurement, the residual
scaling as $u^{2.04}$.  On $64\times64$ and both rectangular orientations it is within
$0.3\%$ at $\ell=10^{-1}$.

\subsection{Inside the loop, and outside the fitted spectra}
\label{sec:loop}

\emph{The predictions survive the loop.}  Over 108 coupled runs (six precisions, three
stopping policies, three seeds, two windows) none diverged, and
$\mathrm{Spearman}(\log\text{pred},\log\text{orth}) = +0.93$ against $-0.05$ for best
validation loss: the law predicts the arithmetic, not the loss.

\emph{The constant is a calibration (held-out shapes and spectra).}  On the logspace control the ratio of \eqref{eq:cert} to measurement
is $0.43$--$0.97$; on a two-plateau spectrum the same constant under-states the error by
$4$--$5\times$ (\S\ref{sec:limitations}).  The first-order formula transfers where the
constant does not ($0.51$--$0.98$, median $0.89$), and the ratchet is identical in all four
families at fixed $\ell$: the law's domain is a spectrum family as much as a precision.
\emph{A dimension-invariant form.}  Dividing the spectrum-only proxy
$\hat v\norm{W}_F$ ($W=[1/(\sigma_i+\sigma_j)]$, $\hat v=\norm{\hat M}_F/(\sqrt3 d)$) by
$\sqrt d$ leaves one number: $6.07\,u\,\hat v\norm{W}_F/\sqrt d$ at the median over 240
cells (CI $[5.98,6.14]$), $89\%$ inside $\pm25\%$, $6.04$ on 64 held-out cells
(Table~\ref{tab:dsinv}).  Its two factors are \S\ref{sec:profile}: a derived $2\sqrt2$
and the spectrum's own $g$.

\subsection{The participation factor, closed}
\label{sec:profile}

For a state of spectrum $s$ let $D_{ij}=1/(s_i+s_j)$ and, for an exchangeable $F$ with
unit entry rms, $g(s)=\norm{(F-F^\top)\circ D}_2/(2\sqrt2\,\norm{D}_F/\sqrt d)$; $g=1$ is
the exchangeable prediction.  Over 180 cells (four families, three $\ell$,
$d=64$--$1024$, three seeds) $g$ is $1.31$ at the median (range $0.98$--$2.10$), and it is
the edge of a variance-profile Wigner operator: coarse-grain $V=2D^2$ into $K$ contiguous
blocks of equal count, $B_{kl}=\overline{V}_{kl}\sqrt{m_km_l}$, read
$2\sqrt{\lambda_{\max}(B)}$.
\footnote{Forming $D$ and $B$ is $O(d^2)$ arithmetic on $d^2$ numbers, so $g$ is a
quantity computed from the spectrum, not a matrix-free sensor: the bound that needs no
spectrum is \eqref{eq:bound}.}
At $K=1$ this returns $2\sqrt2$ exactly; at $K=4$ it lands
within $1.2\%$ of $g$ at the median ($97\%$ of 120 cells inside $\pm10\%$; at
$K=d$, $3.6\%$); a participation-ratio guess misses by $11\times$.  The earlier negative
was this operator with a $\sqrt d$ misplaced.  It follows the decades a spectrum spans,
not $d$: at fixed (family, $\ell$), $g(1024)/g(64)$ spans $0.87$--$1.23$ (median $1.06$),
so the earlier drift compared different $\ell$.  It does not follow the loop: giving each
of the $13$ roundings the one-step law of its state over-predicts the measured floor by
$1.27$--$2.36\times$, while the quadrature of the $13$ linearised responses matches it to $0.92$--$1.12$ (median $1.03$) over 96 cells that
reach $d=512$ and $97\%$ of them within $\pm10\%$; the responses have mean $|\cos|$ $0.02$
and $84\%$ of the sum is on the input cast.  The multiplicity is counted, not modelled
(Appendix~\ref{app:constant}).

\subsection{Real training: a flat loss and a sharp boundary}
\label{sec:realtrain}

To see whether the frozen-trajectory picture survives a loop that reacts, we run the
kernel inside training on \texttt{digits} (a five-layer MLP $64\to128^4\to10$, five seeds,
and a two-block residual MLP, three) and on \texttt{diabetes} (a four-layer regression MLP,
three), Muon on every hidden weight (seven conditions, Table~\ref{tab:realtrain}).

\emph{The loss does not move.}  On \texttt{digits} the training loss spans
$5.77$--$7.39\times10^{-5}$ and the accuracy $0.977$--$0.980$ over all seven
conditions; on \texttt{diabetes} the MSE spans $1.30$--$1.57\times10^{-3}$, and at
twelve bits (E5M2, $K=2$) an update error of $0.73$ still costs $1.65\times$ the fp32 MSE
(Appendix~\ref{app:constant}).

\emph{What does move is stability.}  At the shipped safety E5M2 blows up in $6$ of $11$
seed-runs within $400$ steps (digits $2/5$, residual $1/3$, diabetes $3/3$), E4M3 once, and
nothing else; over $1{,}200$ steps and $25$ seeds every E5M2 run at $1.01$ fails in both
tasks, as does $17$ of $25$ at $1.035$ on \texttt{digits} (Table~\ref{tab:census}).  The
format is not quantising badly (flushed values have median $0$ and E5M2's relative error
is $5.3\%$, what four bits imply): the mode is dynamical and completes inside one call
(\S\ref{sec:escape}).

\subsection{Why a run escapes: one call, one window edge}
\label{sec:escape}

The failure is not a training instability: in each failing run the weights
move by $1$--$2\%$ per step and then, inside one call, the norm jumps by twenty
orders of magnitude (diabetes, seed 0, step 24: $\norm{O}_F=4.9\times10^{44}$).  The
state's top within that call traces the escape (Fig.~\ref{fig:escape}a): the exact orbit
lands on the certified edge of each of the six steps, the E5M2 orbit
leaves it at step 2, and all four instrumented runs repeat the trace.  A rounding of a
tenth of a percent at step 1 lands the state past the quintic's extremum, step 2 returns it
$14\%$ above its edge instead of below, and beyond the edge the
gain is $2.2\times$ at $1.5$ times the edge and $9.3\times$ at twice.

Two counterfactuals change one thing: rescaling the state into its window whenever a step
would leave it makes all $50$ long runs finish, while a $10\%$ cut in the learning rate
leaves $7$ of $11$ failing.  The mode is the escape, not the step
size; the curvature loop is excluded directly: with
$G=\lambda_{\max}(\mathrm{diag}(\mathrm{step})^{1/2}H\,\mathrm{diag}(\mathrm{step})^{1/2})$
from power iteration on Hessian-vector products, $G>1$ on no trajectory (median
$\le0.26$, max $0.83$ in five of six).  What the safety buys is the excess a window
tolerates: $0.0075$ at $1.01$ and $0.048$ at $1.035$ for four bits, the critical safety
$1+c\,u$ of \S\ref{sec:fp8} with $c\approx0.5$ (Fig.~\ref{fig:escape}b), for $0.4\%$ of
the update's Frobenius norm.

\emph{The risk is per call, and it grows.}  Over $1{,}200$ steps and $25$ seeds
(Table~\ref{tab:census}) one call at the shipped safety fails at $1$ in $1{,}600$ on
\texttt{digits} and $1$ in $109$ on \texttt{diabetes} ($14.7\times$, call-level exact
$p=3.5\cdot10^{-17}$), while both fail $25$ of $25$ runs: the run count cannot show the
difference.  The hazard rises --- on \texttt{digits} from $1$ in $6{,}400$ in the first two
hundred steps to $1$ in $370$ in the fourth and fifth hundreds.  At $1.035$ it falls to $1$
in $5{,}900$ on \texttt{digits} and to none in $90{,}000$ calls on \texttt{diabetes}; the
clamp leaves none in $210{,}000$ calls of either.

\begin{figure}[t]
\centering
% --- BEGIN figures/fig2-escape-phase.tex ---
% One call of the iteration that escapes, and the excess a window tolerates.
% Coordinates: figures/data/fig2-*.dat, written by data/make_figure_data.py from
% pilot23_results.json (e_escape_traces.diabetes_s0_e5m2: exact, tops, u; w_windows).
\begin{tikzpicture}
\begin{groupplot}[
  group style={group size=2 by 1, horizontal sep=1.7cm},
  width=0.47\linewidth, height=4.3cm,
  tick label style={font=\footnotesize},
  label style={font=\footnotesize},
  title style={font=\footnotesize},
  legend style={font=\scriptsize, draw=none, fill=none},
  axis line style={thick}, grid=major, grid style={densely dotted, black!18},
]
\nextgroupplot[
  ymode=log,
  xlabel={step of the call}, ylabel={largest singular value},
  xmin=-0.3, xmax=6.3, xtick={0,1,2,3,4,5,6},
  ymin=0.5, ymax=1e47,
  title={(a) E5M2, safety $1.01$: one escape},
  legend pos=north west, legend cell align=left,
]
\addplot[black, thick, mark=o, mark size=1.6pt]
  table[x=step, y=exact] {figures/data/fig2-escape.dat};
\addplot[black, thick, dashed, mark=square*, mark size=1.7pt]
  table[x=step, y=e5m2] {figures/data/fig2-escape.dat};
\addplot[gray, thick, densely dotted, mark=triangle*, mark size=1.6pt]
  table[x=step, y=edge] {figures/data/fig2-edge.dat};
\legend{exact arithmetic, E5M2, window edge}

\nextgroupplot[
  xlabel={safety $s$}, ylabel={admissible excess $\delta^\star$},
  xmin=0.995, xmax=1.055, ymin=0, ymax=0.072,
  xtick={1.0, 1.02, 1.035, 1.05},
  title={(b) the start excess a window tolerates},
  legend pos=north west, legend cell align=left,
]
\addplot[black, thick, mark=square*, mark size=1.6pt]
  table[x=safety, y=delta4] {figures/data/fig2-delta.dat};
\addplot[black, thick, dashed, mark=o, mark size=1.6pt]
  table[x=safety, y=delta5] {figures/data/fig2-delta.dat};
\addplot[gray, thick, dotted, mark=triangle*, mark size=1.7pt]
  table[x=safety, y=delta8] {figures/data/fig2-delta.dat};
\legend{E5M2 ($p=4$), E4M3 ($p=5$), bf16 ($p=8$)}
\end{groupplot}
\end{tikzpicture}
% --- END figures/fig2-escape-phase.tex ---

\caption{Why a call escapes, and what damping buys (\texttt{pilot23}).
(a) \texttt{diabetes}, seed $0$, layer $0$, step $24$: the exact largest singular value
follows the certified upper edge of each step ($1.00$, $1.92$, $1.71$, $1.29$, $1.02$,
$1.00$), the E5M2 orbit leaves it at step $2$ and climbs $44$ decades in four steps.
(b) The largest start excess $1+\delta^\star$ a window tolerates, per step: $0.0075$ at
$s=1.01$ for four bits and $0.048$ at $1.035$, across which the $400$-step census goes from
$6$ failures to none and the $1{,}200$-step hazard on \texttt{digits} from $1$ in $1{,}600$
to $1$ in $5{,}900$ calls.}
\label{fig:escape}
\end{figure}

\subsection{The step budget on real spectra}
\label{sec:steps}

The decisions of \S\ref{sec:alloc} are about where bits go; the one a kernel makes first
is how many steps to run.  Table~\ref{tab:steps} (appendix) sweeps it on 80 momentum
matrices of a five-layer \texttt{digits} run ($128\times128$ and $128\times64$, sampled
every 20 steps) in the design window $\ell_0=10^{-2}$; bits are the input cast plus one
stored width per step.  Fig.~\ref{fig:pareto} is the same sweep read against the budget.

\emph{The step budget.}  In every format the window error stops falling at $K=5$: the
sixth step spends eight bits for under $1.5\%$, and the smallest count within $1.15\times$
of its best error is $K=5$ for $69$ of $80$ matrices and $K=4$ for the rest.  Downgrading a slot is never the
cheaper trade: E5M2 on one slot of a six-step bf16 run multiplies the error by $5.5$ at
the first step and $12.6$ at the second, and dropping the sixth step saves eight bits for
nothing.

\emph{Mixing precision is dominated.}  The $40$-bit hybrid errs by $1.45\times10^{-1}$,
six times worse than bf16 at four steps for the same bits ($2.43\times10^{-2}$), and its
E4M3 body ($44$ bits, $7.36\times10^{-2}$) is still three times worse; a five-step E4M3
update spends $30$ bits for $1.24\times10^{-1}$.  The frontier is (E5M2, $K=4$), (E4M3,
$K=4$--$5$), (bf16, $K=4$--$5$): the $40$-bit and $44$-bit hybrids are on none of it
(Fig.~\ref{fig:pareto}), and
the window where they are not dominated is in \S\ref{sec:fp8}.

\begin{figure}[t]
\centering
% --- BEGIN figures/fig3-pareto-frontier.tex ---
% The measured bit budget against the window error, on the 80 momentum matrices.
% Coordinates: figures/data/fig3-*.dat, written by data/make_figure_data.py from
% pilot21_results.json (D.ksweep: the three formats by K; D.ladder: the mixed rule).
\begin{tikzpicture}
\begin{axis}[
  width=0.72\linewidth, height=5.4cm,
  tick label style={font=\footnotesize},
  label style={font=\footnotesize},
  legend style={font=\scriptsize, draw=none, fill=none,
    at={(1.02,1.0)}, anchor=north west, legend cell align=left},
  axis line style={thick}, grid=major, grid style={densely dotted, black!18},
  ymode=log,
  xlabel={bits (input cast $+$ one width per step)},
  ylabel={window error, median over 80 matrices},
  xmin=8, xmax=62, ymin=8e-3, ymax=1.1,
  xtick={10,20,30,40,50,60},
]
\addplot[black, thick, mark=square*, mark size=1.6pt]
  table[x=bits, y=err] {figures/data/fig3-bf16.dat};
\addplot[black, thick, dashed, mark=o, mark size=1.6pt]
  table[x=bits, y=err] {figures/data/fig3-e4m3.dat};
\addplot[black, thick, densely dotted, mark=triangle*, mark size=1.7pt]
  table[x=bits, y=err] {figures/data/fig3-e5m2.dat};
\addplot[gray, thick, mark=diamond*, mark size=2.2pt, only marks]
  table[x=bits, y=err] {figures/data/fig3-ladder.dat};
\addplot[gray!60, very thick] table[x=bits, y=err] {figures/data/fig3-frontier.dat};
\node[font=\scriptsize, anchor=south west] at (axis cs:25,0.128) {E4M3, $K=4$};
\node[font=\scriptsize, anchor=south west] at (axis cs:40,0.023) {bf16, $K=4$};
\legend{bf16, E4M3, E5M2, mixed rule, measured frontier}
\end{axis}
\end{tikzpicture}
% --- END figures/fig3-pareto-frontier.tex ---

\caption{The bit budget against the window error on the 80 momentum matrices of
Table~\ref{tab:steps} (\texttt{pilot21}).  Lines are the three formats at $K=2$--$6$;
diamonds are the mixed rule at $K=6$ (seven E5M2-body and seven E4M3-body splits, the
all-bf16 end shared); the grey staircase is the frontier of every measured point.  It bends
at two knees --- (E4M3, $25$ bits) and (bf16, $40$) --- where the error falls more than
fivefold for eight bits, and between them it is almost flat: (E4M3, $K=5$) at $30$ bits and
the all-E4M3 end of the ladder at $38$ are on it for under $3\%$ each, with the $35$-bit
E4M3 point above the $30$-bit one and the ladder's $41$- and $44$-bit points dominated by
(bf16, $40$).  Every split that mixes the body's format lies above the staircase.}
\label{fig:pareto}
\end{figure}

\subsection{The exponent range of fp8}
\label{sec:fp8}

fp8 does not break the law where the law is used: in the normal range each format's grid
\emph{is} the kernel's (E4M3 to $q(\cdot,5)$ on $4000$ of $4000$ window points, E5M2 to
$q(\cdot,4)$).  At $d=64$, $K=6$, five seeds (Table~\ref{tab:fp8}) the
E4M3 floor is within $0.3$--$3.2\%$ of the matched unlimited-exponent floor at four
dynamic ranges and E5M2 is $1.86$--$2.09\times$ larger, the factor $2$ one significand
bit predicts.

\emph{The first slot is where precision pays --- inside the window.}  At four bits and
$\ell=10^{-3}$ the shipped $1.01$ blows up in 12 of 12 runs; the smallest damping that
stops it, $1.0313$, moves the benign-range floor by at most $4\%$ and yet leaves
the $\ell=10^{-3}$ floor at the cap in 12 of 12.
Spending the first slot works: cast and step one at bf16 with the other five at $p=4$
takes $\ell=10^{-2}$ from $4$ to $0$ blow-ups and the floor from $0.877$ to $0.243$
($3.6\times$), and two bf16 steps give $0.146$ ($6.0\times$) at $5.33$ bits per
slot against $0.411$ for all-E4M3 at five (Table~\ref{tab:e5m2}).  This is the one regime
where the hybrid wins; under a training kernel's normalisation the same rule is
six times worse than bf16.

\subsection{What to do in the next Muon step}
\label{sec:rules}

Table~\ref{tab:rules} prices the four decisions this section supports.
% --- END sections/06-experiments.tex ---

% --- BEGIN sections/07-related.tex ---
\section{Related Work}
\label{sec:related}

\paragraph{Orthogonalized updates.}
Muon replaces the momentum matrix by its polar factor before the weight
update~\cite{liu2025muon}; that factor is computed by a fixed Newton--Schulz-type
iteration whose convergence and error analysis are
classical~\cite{higham1986polar,higham2008functions,golub2013matrix}, and
Shampoo~\cite{gupta2018shampoo} is the other line of preconditioned tensor optimization.
Amsel et al.~\cite{amsel2025polar} derive minimax quintic coefficients for exact
arithmetic and add the prescriptions this paper takes as its object; the empty design
axis (Appendix~\ref{sec:design}) is consistent with that optimality holding only in
exact arithmetic.  Du and Su~\cite{du2026polarswitch} put rounding into the design as a constraint; we
bound what such a design can buy.
Huang~\cite{huang2026orthogonalization} reports a non-monotone relation between
orthogonalization accuracy and training quality, which the two-term law explains in shape
if not in constant.

\paragraph{Low precision and bit allocation.}
The formats are standardized~\cite{micikevicius2022fp8} and mixed-precision training is
standard~\cite{micikevicius2018mixed}; optimizer states have been quantized
blockwise~\cite{dettmers2022optimizers}, and limited precision with stochastic rounding
has been studied for training~\cite{gupta2015limited}.  That line asks what precision the \emph{parameters and states} need; we ask what the
\emph{working matrices of one iteration} need, and the answer is a floor proportional to $u$
with a computable constant and a separate exponent-range test.  Quantization work
allocates bits across layers, channels or tensors~\cite{hariri2025quantize,dettmers2022optimizers};
here allocation is inside one step, across $K+1$ slots, and it is a maximum rather than a
sum, which makes a closed form possible.

\paragraph{Perturbation of the polar factor, and stopping.}
First-order perturbation of the SVD and of orthonormal bases is
classical~\cite{stewart1990perturbation,golub2013matrix}; the two-term form
\eqref{eq:firstorder}, with the projected term written separately because it survives for
rectangular matrices, is what we could not find stated as a working estimate for a
rounded momentum matrix.  Early-exit rules in iterative solvers test a residual the algorithm can see; here the law
runs before the iteration, against an error it cannot see.
% --- END sections/07-related.tex ---

% --- BEGIN sections/08-limitations.tex ---
\section{Limitations}
\label{sec:limitations}

\emph{The kernel is a model, and its constant is a calibration.}  We round the state and
the intermediate to $p$ significant bits and accumulate in float64; fp8's exponent range
is priced separately (\S\ref{sec:fp8}); fused multiply--add, accumulation order and
subnormals in the wide accumulator are not modelled.  The three alternative kernels of
Appendix~\ref{sec:variants} keep the $u$-scale and the $K$-dependence but move $c$ by up
to $2.2\times$.

\emph{The constant moves with the spectrum family.}  The 48 held-out cells of
\S\ref{sec:loop} put the law at $0.43$--$0.97$ of measurement on logspace spectra and at
$4$--$5\times$ optimism on a two-plateau one.  The ratchet is not what moves: at fixed $\ell$ the
prediction is identical in all four families, and a first-order bookkeeping of every
rounding closes these measurements to $2\%$ ($63$--$91\%$ of the response on the input
cast for logspace spectra, $88$--$99\%$ on two-plateau ones).  The normalised pair-weight
operator is one constant across families and dimensions only to a factor, and the profile
factor that closes one spectrum (\S\ref{sec:profile}) does not follow the loop's
injections.  On the rectangular momenta of the Muon loop the same constant under-states the floor by
a median $2.9\times$ ($\hat c=0.41$) because the projected term of \eqref{eq:firstorder} is
not in the proxy, which is why the floor claim on that shape rests on
\eqref{eq:firstorder}.

\emph{The safety factor is a margin measured on a window, not a convergence knob.}
What damping buys is the excess a state may exceed its edge and still converge, $0.008$ at
$s=1.01$ and $0.048$ at $1.035$ for four bits (\S\ref{sec:escape}) --- the measured
critical safety $1+c\,2^{-p}$ with $c=0.04$--$0.50$ (Table~\ref{tab:safety}), $c\approx0.5$
on the tight steps, $22$--$55\times$ the margin bf16 needs.  At $\ell=10^{-3}$ four bits
stay at the cap anyway: stop earlier (\S\ref{sec:steps}) or spend the first slot's
precision where the window is wide (\S\ref{sec:fp8}).

\emph{The step budget is measured on one window.}  The frontier of \S\ref{sec:steps} comes
from the $128\times128$ and $128\times64$ momenta of a five-layer MLP at
$\ell_0=10^{-2}$; the plateau sits where the exact error falls below the floor, so a wider
window moves it.  On the full spectrum every design errs by $1.00$--$1.65$, because the
median dynamic range is $2\cdot10^{-10}$ and the tail is not reached in six steps: we do
not claim the truncated tail is harmless, only that no precision repairs it.

\emph{One scale of model, a flat loss, and no online rule.}  The training measurements stop
at five layers and 64 input dimensions, eleven seed-runs per condition at $400$ steps and
twenty-five at $1{,}200$, and call a
divergence a weight exceeding $10^{3}$: a transformer's momenta may be narrower, which
moves both the profile factor and the plateau.  At this scale the update
error spans $1.6\times10^{-2}$ to $0.73$ while the training loss moves by $1.65\times$;
when $6\%$ would not is a GPU-scale experiment we have not run.  The per-call hazard of
\S\ref{sec:escape} rises with the step count, so $1$ in $1{,}600$ calls is not the rate of a
$10^{9}$-token run; read as a first-order rate it already says the shipped safety cannot
outlast one.
Allocation is verified
for two operands and for per-slot schedules separately, not jointly, and the schedule has
no online form (\S\ref{sec:online}).
% --- END sections/08-limitations.tex ---

% --- BEGIN sections/09-conclusion.tex ---
\section{Conclusion}
\label{sec:conclusion}

The coefficients of a minimax polynomial iteration are chosen for exact arithmetic, and
the iteration does not run in it.  Storing the state and the intermediate in $p$ bits adds
a floor $c\,u\,S_{\mathrm{pert}}(K)$, $u=2^{-p}$; with the exact error it predicts the
measured error to a median ratio of $1.10$ over 90 cells, and survives bf16, fp8's
exponent range and three other kernels up to the constant.

The floor is not a balance of contraction and reinjection but a subspace error: the exact
iteration preserves the polar factor, so no step count, design or schedule can undo an
error created at the first quantisation.  A first-order formula gives that error its size: for a square state the
pair weights $1/(\sigma_i+\sigma_j)$ say which roundings are expensive, and for the
rectangular state a Muon update actually is, the rotation of the column space dominates,
without which the prediction is off by $2.8\times$ in all 500 real momentum matrices.
The one constant left is a geometric factor of the spectrum, which a four-block
variance-profile operator reproduces to $1.2\%$ from the spectrum alone.

Three consequences follow, each measured at equal cost.  A target accuracy is met by the
first $K$ whose law value is below it, with no overshoot beyond $1.16\times$ in 62 cells,
while the crossing reading of the same law loses up to three orders of magnitude.  Bits go
where the max of the per-slot gains puts them, a median $3.06\times$ where the dynamic
range is wide.  Yet on real momentum the decision that pays is earlier: the fifth
step is a plateau, the sixth spends eight bits for under $1.5\%$, and at equal bits a
four-step bf16 update beats the mixed-precision rule sixfold, an advantage of the
$\sigma_{\max}$-anchored window rather than of the arithmetic.  Inside training the loss
registers none of this --- the update error spans $1.6\times10^{-2}$ to $0.73$ while the
loss moves by $1.65\times$ --- and what separates runs is survival: the state's top leaves
its certified window inside one call, and a safety factor of $1.035$ widens that tolerance
sixfold and lowers the hazard of a call $4$--$200\times$.  That is margin, not immunity:
at $1$ in $1{,}600$ calls the shipped safety cannot outlast a long run, and only keeping the
state inside its window left no failure in $210{,}000$ calls.

What is missing is a tight matrix-free certificate, an online rule for the budget, and a
demonstration at scale.  The claim is narrower than a new algorithm: the precision floor of an orthogonalization
step is predictable from the coefficients before the step runs, and the daily decisions
--- how many steps, where the bits go, whether to mix them, how much damping --- follow
from it.
% --- END sections/09-conclusion.tex ---

% --- BEGIN sections/99-appendix.tex ---
\appendix

\section{Reproduction}
\label{app:repro}

Each experiment is one script that writes one JSON file; the text quotes the JSON.
All scripts are pure \texttt{numpy}/\texttt{scipy}, run on a CPU, and the whole
sequence takes about 70 minutes.  The scripts are named after the order in which
they were run, not after the paper's sections.

\begin{table}[h]
\centering
\caption{Measurement scripts.  ``cell'' counts are the number of distinct
(matrix, precision, step count) measurements the script performs, which is the
number the text quotes.}
\label{tab:scripts}
\footnotesize
\setlength{\tabcolsep}{4pt}
\begin{tabular}{>{\raggedright\arraybackslash}p{5.6cm}
                >{\raggedright\arraybackslash}p{6.3cm}
                >{\raggedright\arraybackslash}p{3.4cm}}
\toprule
script & what it measures & used in \\
\midrule
\texttt{pilot1\_polar\_precision} & exact vs p-bit error of NS5/PE/LSQ5 designs & \S\ref{sec:exp} \\
\texttt{pilot1c\_precision\_scaling} & error against $K$ at fixed $p$ (flatness) & \S\ref{sec:exp} \\
\texttt{pilot4\_precision\_placement} & two-operand table, equal-cost allocation & \S\ref{sec:alloc} \\
\texttt{pilot4b\_k1\_separation} & the same table for $K=1,2,3,6$ (quintic) & \S\ref{sec:alloc} \\
\texttt{pilot4c\_k1\_highdeg} & $K=1$ with a degree-17 step ($E_1$ small) & \S\ref{sec:firstorder}, \S\ref{sec:alloc} \\
\texttt{pilot5\_lower\_bound} & single-step LP frontier, degree scaling, $\alpha$ & Appendix~\ref{app:refute} \\
\texttt{pilot5b\_scaling} & $S_{\mathrm{peak}}(\ell)$, $K^\star(\ell)$, plateau constant & \S\ref{sec:exp} \\
\texttt{pilot5c\_composition\_design} & non-monotone design path, $S_{\min}(\varepsilon,d)$ & Appendix~\ref{app:refute} \\
\texttt{pilot6\_s\_pert} & 192-cell competition of certificate definitions & \S\ref{sec:prelim}, \S\ref{sec:cert} \\
\texttt{pilot6b\_floor\_constant} & injected rounding size, $d$-scaling, $u$-linearity & \S\ref{sec:cert} \\
\texttt{pilot7\_pert\_aware\_design} & 37 designs measured, objective argmin, divergence census & \S\ref{sec:design} \\
\texttt{pilot7c\_table1} & certificate table over $p=4..24$, 90-cell stopping & \S\ref{sec:cert} \\
\texttt{pilot8\_muon\_mlp} & coupled Muon runs, frozen replay, rules, dose-response & \S\ref{sec:loop} \\
\texttt{pilot8b\_equilibrium} & subspace/value split, tail extension & \S\ref{sec:invariance} \\
\texttt{pilot10\_subspace\_theory} & invariance, first-order formula, $C(\ell)$, SVD-free bound & \S\ref{sec:cert} \\
\texttt{pilot10b\_first\_step\_schedule} & slot profile, 62-cell stopping, 48-cell schedule & \S\ref{sec:stopping}, \S\ref{sec:alloc} \\
\texttt{pilot11\_momentum\_subspace} & 500 real momentum matrices: two-term formula, $C\ell$ & \S\ref{sec:realmomentum} \\
\texttt{pilot11b\_application\_order} & fixed design and matrix, $K=1,2,3,4$ & \S\ref{sec:alloc} \\
\texttt{pilot12\_arithmetic\_variants} & three other kernels, self-calibration of $c$, plan cost & \S\ref{sec:variants} \\
\texttt{pilot13\_ell\_sensitivity} & design and stopping rule under $\hat\ell=\alpha\ell$, 252 decisions & \S\ref{sec:ellsens} \\
\texttt{pilot14\_real\_allocation} & graded schedule on 500 real momenta, 712 matrix-cells, $\ell$ sweep & \S\ref{sec:premium} \\
\texttt{pilot15\_fp8\_kernel} & E4M3/E5M2 with subnormals and saturation, 5 seeds, 4 dynamic ranges & \S\ref{sec:fp8} \\
\texttt{pilot16\_stats\_and\_lp} & bootstrap intervals, LP grid convergence, integer DP, held-out shapes & \S\ref{sec:loop}, \S\ref{sec:alloc} \\
\texttt{pilot17\_theory\_refinement} & family dependence of the constant, pair recursion, safety boundary & \S\ref{sec:limitations} \\
\texttt{pilot18\_c\_scaling\_and\_e5m2\_fix} & dimension normalisation of the floor constant, four-bit repair sweep & \S\ref{sec:loop}, \S\ref{sec:fp8} \\
\texttt{pilot18b\_bootstrap} & intervals and paired gains on the two above & \S\ref{sec:loop}, \S\ref{sec:fp8} \\
\texttt{pilot19\_large\_d\_and\_decomp} & dimension extrapolation to $d=512$; the constant's factors & \S\ref{sec:loop} \\
\texttt{pilot20\_analytic\_and\_d1024} & $2\sqrt2$ from the theorem-3 operator; $d=768,1024$; the practical rules & \S\ref{sec:loop}, \S\ref{sec:rules} \\
\texttt{pilot21\_loss\_phase\_and\_scheduler} & the kernel inside training (digits, diabetes, 11 seed-runs); step budget on 80 real momenta; online sensors & \S\ref{sec:realtrain}, \S\ref{sec:steps}, \S\ref{sec:online} \\
\texttt{pilot21b\_kmin} & per-matrix step counts on the same 80 momenta & \S\ref{sec:steps} \\
\texttt{pilot22\_participation\_factor} & the profile factor $g$ on 180 cells; the block-operator closure; per-injection coefficients & \S\ref{sec:profile} \\
\texttt{pilot23\_divergence\_mechanism} & the window and the escape; the four counterfactuals; the loop gain & \S\ref{sec:escape} \\
\texttt{pilot24\_large\_d\_multiplicity} & the multiplicity to $d=512$; per-injection quartiles; the ratchet's bound & \S\ref{sec:profile}, Appendix~\ref{app:bound} \\
\texttt{pilot25\_long\_census} & 1,200 steps $\times$ 25 seeds: the hazard of a call, with intervals and tests & \S\ref{sec:realtrain}, \S\ref{sec:escape}, Table~\ref{tab:census} \\
\texttt{make\_figure\_data} & every plotted coordinate, read from the JSON files above & Figs.~\ref{fig:ratchet}--\ref{fig:pareto} \\
\bottomrule
\end{tabular}
\end{table}

\paragraph{Determinism.}
The momentum sequence used in \S\ref{sec:realmomentum} is reproduced bit for bit
by rerunning the exact-arithmetic trajectory with the same seed; the replay was
checked against the sequence of the run that produced \S\ref{sec:loop} (median
dynamic range $0.1499$, best validation loss $0.6534$, accuracy $0.762$, all
identical).  The matrix of \S\ref{sec:alloc} is reproduced by seeding
\texttt{numpy} as in the single-step script; its exact error $E_1$ matches to
seven digits.

\paragraph{Design generator.}
The coefficient generator (a Newton iteration on the fifth-degree equioscillation point)
returns non-finite coefficients at five of three hundred grid values of $\ell$, from
$3.2\times10^{-2}$ to $0.46$.  The held-out grid that would have used $3\times10^{-2}$
tests the design and raises instead of running the kernel with a non-finite one, and uses
$5\times10^{-2}$ there.

\paragraph{Archive.}
The scripts, their JSON results and logs, and the validation script are distributed as
\path{iclr-2028-artifact.tar.gz}; \texttt{MANIFEST.sha256} pins each experiment file
and hashes to \texttt{5aebe2b2772e38cc}\allowbreak\texttt{e2c054c2860d32c3}\allowbreak\texttt{a71845610635d1e5}\allowbreak\texttt{1e2ad70c579f2375}, so the numbers in the tables can be checked
against the code that produced them without re-running it.

\section{The Negative Results in Full}
\label{app:refute}

\paragraph{No lower bound of the form $S\ge c\,\varepsilon^{-\alpha}$.}
Three measurements, all from linear programs over odd polynomials of degree
$d\in\{3,5,7,9,13,17\}$ with the constraint $\max_{[\ell,1]}|P-1|=\varepsilon$ and
the objective the coefficient $\ell_1$ norm.

(i) At fixed error the minimal sensitivity decreases monotonically in the degree:
\[ S_{\min} = 1.178,\ 0.937,\ 0.648,\ 0.371,\ 0.114,\ 0.0363
   \quad \text{for } d = 3,5,7,9,13,17, \]
at $\ell=0.3$, $\varepsilon=0.6$ --- a factor $32$ across the
family.  So small $\varepsilon$ and small $S$ are simultaneously achievable.

(ii) The local exponent of $S_{\min}(\varepsilon)$ is not a constant, and it depends
on $\ell$ by two orders of magnitude:

\begin{table}[h]
\centering
\caption{Local exponent $\alpha$ of $S_{\min}\propto\varepsilon^{-\alpha}$ at three
dynamic ranges, from the linear program (log--log fit over the admissible range of
$\varepsilon$; fit residual $0.19$--$0.26$ in $\log$).}
\label{tab:alpha}
\begin{tabular}{rcccccc}
\toprule
$\ell$ & $d=3$ & $d=5$ & $d=7$ & $d=9$ & $d=13$ & $d=17$ \\
\midrule
$0.3$   & 1.43 & 0.81 & 0.56 & 0.45 & 0.32 & 0.22 \\
$0.1$   & 7.50 & 3.55 & 2.39 & 1.78 & 1.13 & 0.77 \\
$0.03$  & 21.4 & 13.1 & 7.11 & 5.02 & 4.60 & 3.13 \\
\bottomrule
\end{tabular}
\end{table}

(iii) Along a design path at fixed $(\ell,K,d)$ the error falls by $500\times$ while
the sensitivity is non-monotone:
\[ \varepsilon = 0.843,\ 0.782,\ 0.705,\ 0.612,\ 0.505,\ 0.386,\ 0.263,\ 0.150,\ 0.061,\ 0.021,\ 0.0016, \]
\[ S = 155.6,\ 215.0,\ 286.4,\ 366.7,\ 448.5,\ 516.8,\ 548.8,\ 517.1,\ 404.3,\ 228.0,\ 60.4. \]
The local slope
$\mathrm{d}\log S/\mathrm{d}\log\varepsilon$ runs from $-4.26$ to $+0.53$.

\paragraph{Why.}
The work of the lower bound would have to happen in the boundary layer
$[0,\ell]$, where the polynomial has to leave its value at the origin; the
sensitivity is measured on $[\ell,1]$ only, and raising the degree thins the
boundary layer without paying on the interval.  At $\ell=10^{-3}$ a single step
cannot achieve an error below $0.873$ even at degree $65$, so the regime the paper
works in is one where iteration is mandatory and no single-step bound can be tight.

\paragraph{Two operands.}
Table~\ref{tab:k1} is the two-operand question across the state and the intermediate.
The whole column of $B$ from 6 to 24 bits is worth $1.12\times$ at $p_X=8$ and
$71.7\times$ at $p_X=16$: the bottleneck is whichever operand is coarser, and at equal
total bits $(16,16)$ gives $\e{2.14}{-3}$ against $\e{2.90}{-2}$ for $(24,8)$.  The
operands are interchangeable to $1\%$ at $K=1$ and separate by $2.5\times$ at $K=2$, and
the quintic keeps that separation at $K=6$.  On the real spectra it is larger: $(24,8)$
gives $\e{5.74}{-2}$ against $\e{2.44}{-4}$ for $(16,16)$, a factor $236$.

\paragraph{The design axis is empty.}
\label{sec:design}
The law invites one more use: minimize $J(K)$ over the coefficients.  Three measurements
of a $37$-design family (cushion, safety, first-step damping $\tau$; $d=64$) say no.
(i) \emph{No interior optimum}: the argmin over
$(\text{cushion},\text{safety},\tau,K)$ sits on the scan boundary ($K=12$, $\tau=2.2$)
in every cell, because a step costs arithmetic and $J$ does not count it.  (ii) \emph{The
last step's derivative is not a handle}: with $p'(1)=\lambda$ fixed, $J$ is minimized at
$\lambda^\star=-0.0353$, the shipped step's own value, while $S_{\mathrm{pert}}$ is
identical ($609.6$) for all $|\lambda|\le0.5$.  (iii) \emph{Measured, the shipped design
is already there}: at $\ell=10^{-3}$, $K=8$, bf16 the $35$ non-divergent designs span
$0.232$--$0.768$, the shipped point is $2.578\times10^{-1}$ ($1.112\times$ from the best),
and the objective's argmin is $1.58\times$ \emph{worse}.  The reason is in
\eqref{eq:firstorder}: the subspace gain is set by the whole spectrum while
$S_{\mathrm{pert}}$ is a scalar of the one-dimensional recurrence, so $c$ is not
design-independent --- over the 37 designs it ranges over $-0.08$--$0.25$ with
$\mathrm{Spearman}(S_{\mathrm{pert}},\mathrm{err})$ between $-0.84$ and $-0.91$.  Designs
with safety $1.0$ diverge at bf16 in six of the 37 cells, so the shipped $1.01$ is a
divergence margin (\S\ref{sec:limitations}).

\section{A Matrix-Free Bound, and Why No Spectral Bound Exists}
\label{app:bound}

\begin{proposition}[Computable upper bound]
\label{prop:bound}
For full-rank $A$, $B$,
$\norm{\polar(A)-\polar(B)}_2 \le 2\norm{A-B}_2/(\sigma_{\min}(A)+\sigma_{\min}(B))$.
With $A = q(\hat M,p)$, $B = \hat M$ and $\sigma_{\min}(A) \ge \ell - \norm{E}_2$,
\begin{equation}
\mathrm{sub}_{\mathrm{err}} \;\le\; \frac{2\norm{E}_2}{2\ell - \norm{E}_2},
\qquad E = q(\hat M,p) - \hat M ,
\label{eq:bound}
\end{equation}
where $\ell$ is the design's assumed dynamic range and $\norm{E}_2$ is obtained
with matrix--vector products only: power iteration on $E^\top E$ reproduces the
SVD value to a relative $1.2\times10^{-15}$ in all 24 bound cells.
\end{proposition}

Equation~\eqref{eq:bound} is valid in all 24 cells tested and loose by
$2.1$--$14.0\times$ (median $5.7\times$); the slack grows with the dynamic range
because $1/(2\ell)$ assumes the worst orientation while rounding has a fixed
participation.  Replacing $\norm{E}_2$ by the exactly computable $\norm{E}_F$ costs
another factor: the Frobenius form is valid but loose by $6.3$--$45\times$ (median
$24\times$).  A matrix-free guarantee therefore exists, at about one order of
magnitude in accuracy; the tight prediction of \eqref{eq:firstorder} is not
matrix-free.

\begin{proposition}[Spectral information is not enough]
\label{prop:nospectral}
No function of the singular values of $X$ upper-bounds
$\mathrm{sub}_{\mathrm{err}}(X) = \norm{\polar(X)-\polar(M)}_2$.
\end{proposition}

Take a random orthogonal $Q$: its singular values are all $1$, so any spectral
functional takes the value it takes at $\polar(M)$, while
$\norm{\polar(Q)-\polar(M)}_2 = 2.000$ (measured: $\norm{QQ^\top-I}_2 =
1.6\times10^{-15}$).  The subspace is free once the spectrum is fixed.  On our grid
the two quantities co-vary ($\mathrm{Spearman}(\norm{XX^\top-I},
\mathrm{sub}_{\mathrm{err}}) = +0.83$, 32 cells), but only because both grow with
the dynamic range: a correlation is not a certificate.  This is why
\eqref{eq:cert} contains the orbit sensitivity of the \emph{design}, and why the
floor analysis contains $\norm{E}$, the size of the start error, rather than
anything measured from $X$ alone.

\paragraph{The ratchet, and the channel it certifies.}
Definition~\ref{def:ratchet} is a maximum, and it is worth stating exactly what it
bounds.  Write the $p$-bit scalar orbit as $\hat x_t = p_t(\hat x_{t-1}) + \delta_t$ with
$|\delta_t|\le\rho u$ --- one cast and, per step, the rounding of $B$ and of the state,
$2K+1$ injections in all --- and let $x_t$ be the exact orbit it shadows.

\begin{proposition}[What the ratchet bounds]
\label{prop:ratchet}
Let $p_l\in C^1(W)$ on a window $W$ containing the orbit of every $\sigma\in[\ell,1]$ and
its perturbations, $\rho=\max_W|x|$.  Then, for $u\to0$,
\begin{equation}
|x_K-\hat x_K| \;\le\; (2K+1)\,\rho\,u\,S_{\mathrm{pert}}(K) \;+\; O(u^2),
\label{eq:ratchet}
\end{equation}
and with the off-diagonal gain $\gamma_l(x,y)=a_l+b_l(x^2+y^2)+c_l(x^4+x^2y^2+y^4)$ in
place of $|p_l'|$ in the same prefix maximum, call it $G_{\mathrm{pert}}(K)$, the subspace
error obeys the same inequality.  In neither case may $S_{\mathrm{pert}}(K)$ be replaced by
the terminal $S(K)$.
\end{proposition}

\begin{proof}
$\hat x_t = \phi_t(\hat x_{t-1})+\delta_t$ with $\phi_t$ differentiable on $W$, so the
first-order response to the injection at step $j$ is
$\delta_j\prod_{l>j}p_l'(x_{l-1})$; the response to $\delta_j$ summed over the $2K+1$
injections is bounded by $\rho u\sum_j\prod_{l>j}|p_l'|\le(2K+1)\rho u\max_j\prod_{l>j}|p_l'|$,
and $\prod_{l>j}|p_l'|\le S(K)\le S_{\mathrm{pert}}(K)$ by Definition~\ref{def:ratchet} ---
the suffix product is one of the products $S(K)$ maximises over.  The $O(u^2)$ term is the
second-order remainder of the composition, once per injection, and it can be given its own
coefficient: with $M=\max_W|p_\ell''|$ and $B=\max_W|x|$, the recursion
$e_t=p_t'(x_{t-1})e_{t-1}+\delta_t+\tfrac12p_t''(\xi_t)e_{t-1}^2$ with
$|e_{t-1}|\le(2t-1)BuS_{\mathrm{pert}}(t-1)$ propagates the quadratic parts through the
same prefix products, so
\begin{equation}
|x_K-\hat x_K| \le (2K{+}1)BuS_{\mathrm{pert}}(K)
+ \tfrac12 M B^2 u^2 S_{\mathrm{pert}}(K)\sum_{t=1}^{K}(2t-1)^2 S_{\mathrm{pert}}(t-1)^2
+ o(u^2).
\label{eq:second}
\end{equation}
At the shipped point ($\ell=10^{-3}$, $K=7$, $u=2^{-8}$, $B=1$, $M=2.0\cdot10^3$ over the
visited window) the quadratic coefficient exceeds the first-order term by orders of
magnitude: \eqref{eq:second} is a $u\to0$ statement, and at four bits what certifies
\eqref{eq:ratchet} is the measurement (24 of 24 cells, $16$--$1048\times$ slack), not the
remainder bound.  We claim it as a first-order envelope at four bits.  The pair channel is the
same argument: an injection's off-diagonal component is multiplied, step by step, by
$\gamma_l$ at the pair's own orbit, so the triangle inequality and the definition of
$G_{\mathrm{pert}}$ give it.  The terminal value fails because $S(K)$ maximises over the
injection points of \emph{one} prefix, while the same products regrouped by prefix can
exceed it --- and do when the last step collapses (Fig.~\ref{fig:ratchet}a).
\end{proof}

Two measurements test it (\texttt{pilot24}).  The certified bound holds in all 24 cells
tested ($d=64$, $p=8$, $K=6,7$, four families, three $\ell$) with a slack of
$16$--$1048\times$ (median $41\times$), computed with $\rho=1$: the inequality is true and
loose by about the ratio of $2K+1$ to the calibrated $c$, which is what sign cancellation
does (\S\ref{sec:profile}).  The pair ratchet is not a small correction to the scalar one:
$G_{\mathrm{pert}}/S_{\mathrm{pert}}$ is $0.62$--$1.21$ (below one in 16 of 24 cells), and
the constant it leaves, $\mathrm{err}/(u\,G_{\mathrm{pert}})$, is $0.011$--$0.912$
(median $0.44$) --- below one in every cell.  So the subspace channel needs no constant
above one, and the $0.142$ the text reports is the scalar surrogate of that envelope.  The
two constants are unstable and tight for the same reason.  The scalar inequality replaces
each signed response by its modulus, so the $2K+1$ injections add at full size and the
measured error sits $41\times$ (median) below the bound only through a cancellation that
belongs to the data, not to the design --- which is why what the modulus leaves over moves
over two orders of magnitude.  The pair channel keeps each pair's off-diagonal gain and,
unlike the modulus, does not let the diagonal and off-diagonal parts of an injection be
traded against each other from step to step; that is what makes its envelope hold at one in
every cell (\S\ref{sec:profile}), and its price is a constant that has to be measured
rather than read off $S_{\mathrm{pert}}$.

What no condition on the spectrum or on the coefficients can do is make
$S_{\mathrm{pert}}$ a bound for the \emph{measured} floor.  The condition that would is
$\gamma_l(x,y)\le\max_{u\in[\min(x,y),\max(x,y)]}|p_l'(u)|$ on every reachable pair, and it
fails wherever a step has an interior zero of $p'$: a Newton step has $p'(1)=0$ and
$\gamma(1,1)=1/2$, and the ratio is unbounded as the pair approaches that point.  Over the
reachable pairs and steps the worst ratio is $1.8\times10^3$ at $\ell=10^{-3}$ and
$9\times10^9$ at $\ell=10^{-1}$, so the domination holds in none of the nine designs
tested, and the shipped sequence's own last step at $\ell=10^{-3}$, $K=7$ is
$(1.875,-1.250,0.375)$: $p'(1)=0$, $\gamma(1,1)=0.500$.  That is why the text reports
$S_{\mathrm{pert}}$ with a calibrated constant and names the pair channel as what the
multiplicity counts, instead of claiming a theorem for the floor it measures.

\section{Is the Kernel Model the Arithmetic?}
\label{sec:variants}

A reviewer is entitled to ask whether $c=0.142$ is a constant of our kernel model rather
than of the arithmetic.  We re-ran the shipped design under three other kernels on 12
cells each ($\ell\in\{10^{-2},10^{-3},10^{-4}\}$, $d\in\{64,128\}$, $p\in\{8,11,16\}$)
and report predicted/measured in Table~\ref{tab:variants}.  \emph{float32 accumulation}
gives $0.84$--$1.70$ (median $1.13$) and the same $2.629\times10^{-1}$ at the reference
cell as our kernel, so the wide-accumulator assumption is not what the constant measures;
\emph{eager rounding} raises the floor by up to $2.2\times$ (ratio $0.45$--$1.68$) and
\emph{stochastic rounding} gives $0.56$--$1.70$.  The constant moves by about a factor $2$
across kernels while the $u$-scale and the $K$-dependence do not.  A caller who does not
trust $0.142$ can fit it: $\hat c=0.109$ from 8 cells predicts a held-out grid to
$0.58$--$1.19$ (median $0.85$, CI $[0.64,1.12]$), for the cost of 8 runs; the plan it
enables takes $0.47$\,ms against $2.3$\,ms for the iteration.

\section{The Floor Constant across Families and Dimensions}
\label{app:constant}

\paragraph{The cast factor.}  Measured against the assumed $\rho=u/\sqrt3$ of the input
cast, the kernel's own cast rms is $1.472$ times larger, and for a log-uniform mantissa
$\sqrt{\mathbb{E}[m^{-2}]}=\sqrt{3/(2\ln2)}=1.471$ --- derived, and reproducing the
measured cast error to $0.6\%$ at every $d$ tested.  What remains is the response of an
incoherent perturbation: on the sixteen cells checked in this way the channel factor is
$4.53$, carried $3.8$ by the input cast and $2.3$ by the in-loop roundings in quadrature
($4.43$, $2\%$ low), so the 240-cell constant is $6.07=1.471\times4.13$.  Gaussian noise
of the same rms reproduces the cast channel to $1.00\times$ and a rank-one perturbation of
the same norm is $4\times$ larger, so the incoherence is real and the channel factor is
what stays measured.

\paragraph{Extrapolation.}  Extending the grid to $d=384$ and $512$ (four times the top of
the fit, 96 cells) leaves the constant at $6.48$, $+6.6\%$ with $88\%$ of cells inside
$\pm40\%$, and the trend fitted on $d\le256$ extrapolates to within $11\%$ of them.  Two
more dimension steps move nothing: over 32 cells at $d=768$ and $1024$ (two flat-start
families, two $\ell$, two $p$, two seeds) the median is $6.99$, $1.01\times$ the pooled
median of the same two families at $d=512$, every cell is inside $\pm40\%$ of $6.07$, and
the trend fitted on $d\le512$ extrapolates to $0.93$--$1.00$ of the measurement on all
$16$ predictions.

\paragraph{The participation, and its closure.}  Of the response factor $4.13$ behind
$6.07$, $2.83$ is derived.  On matrix units the first term of \eqref{eq:firstorder} is the
fixed operator $T(e_{ij})=(\sigma_i+\sigma_j)^{-1}(e_{ij}-e_{ji})$, checked against finite
differences over all $992$ units of a $d=32$ cell to a residual of $10^{-6}$, so its
Hilbert--Schmidt norm is $2(\norm{W}_F^2-\sum_iW_{ii}^2)^{1/2}$ without a fitted constant;
an exchangeable perturbation is then answered by $2\sqrt2\,\rho\norm{W}_F/\sqrt d$ --- a
Wigner factor $2$ and a skew factor $\sqrt2$ --- and a flat spectrum realises it ($2.81$
against $2\sqrt2=2.83$ at $d=1024$).  A real spectrum breaks exchangeability, and the
factor that survives, $g$, is measured over 180 cells at $1.31$ (range $0.98$--$2.10$).
It closes as the edge of a coarse-grained variance-profile operator: with $V=2D^2$ and $K$
contiguous blocks of equal index count, $B_{kl}=\overline{V}_{kl}\sqrt{m_km_l}$ and the
prediction $2\sqrt{\lambda_{\max}(B)}$ are within $1.2\%$ of the measurement at the
median for $K=4$ and inside $\pm10\%$ in $97\%$ of 120 cells; the profile's
$\lambda_{\max}$ is $11$--$23\times$ too large and a row degree $76$--$1657\times$, both
of which are this operator with the normalisation misplaced.  The loop is not a second
factor: the one-step formula and the full $K$-step propagation of the same error agree to
$0.95$--$1.07$ (median $0.998$).

\paragraph{What the loop does add, counted.}  The kernel injects one cast and two errors
per step, and their first-order responses add in quadrature to $0.92$--$1.12$ of the
measured floor over 96 cells, with $84\%$ of the total on the input cast.  Giving every
injection the one-step law of the state it enters, however, over-predicts the floor by
$1.27$--$2.36\times$: the per-injection coefficient divided by $2\sqrt2\,g(s_t)$ spans
$0.003$--$2.6$ (median $0.57$), because a rounding injected into an intermediate is not a
rounding of the state.  So $4.13=2.83\times g\times$(the counted multiplicity), with both
factors computable and neither fitted.

\paragraph{The coefficient, in quartiles.}  A range of $0.003$--$2.6$ hides where the mass
is; Table~\ref{tab:w13} splits it.  The input cast obeys the one-step law almost exactly
(median $0.990$, IQR $[0.944,1.011]$ over $16$ injections), the in-loop state roundings are
within $11\%$ of it at the median, and the intermediate's own rounding is the one that
carries almost no first-order response (median $0.054$): the spread the text quotes is one
channel, not noise, which is exactly why the multiplicity is counted rather than modelled
by a per-state law.  The median is stable in $d$ and in $\ell$ (Table~\ref{tab:w13}), and
the coefficient itself does not drift with $d$ where $g$ is left out: over $1248$
injections of the 96 large-$d$ cells the raw $\kappa$ has median $2.11$, $2.21$, $2.25$,
$2.29$ at $d=64,128,256,512$ (a flat spectrum gives $2\sqrt2=2.83$).

\begin{table}[h]
\centering
\caption{The per-injection coefficient of \S\ref{sec:profile} as quartiles
(\texttt{pilot22}, $208$ injections, $p=8$, $K=6$, $d=64,96$): the response of one rounding
divided by $2\sqrt2\,g(s_t)\,\mathrm{rms}(E)\,\mathrm{scale}_t$, so that $1$ means the
one-step law of the state the rounding enters.  ``$X$'' is the rounding of the state,
``$B$'' the rounding of the intermediate; $\ell=10^{-1}$ and $10^{-2}$ are the two windows
of the grid.}
\label{tab:w13}
\begin{tabular}{lrrrrr}
\toprule
group & min & Q1 & median & Q3 & max \\
\midrule
all injections      & $0.003$ & $0.103$ & $0.569$ & $0.931$ & $2.64$ \\
input cast          & $0.878$ & $0.944$ & $0.990$ & $1.011$ & $1.10$ \\
in-loop, $X$        & $0.469$ & $0.715$ & $0.890$ & $0.959$ & $2.64$ \\
in-loop, $B$        & $0.003$ & $0.012$ & $0.054$ & $0.236$ & $2.21$ \\
$\ell=10^{-1}$      & $0.003$ & $0.038$ & $0.696$ & $0.902$ & $1.10$ \\
$\ell=10^{-2}$      & $0.010$ & $0.149$ & $0.539$ & $0.951$ & $2.64$ \\
logspace            & $0.003$ & $0.149$ & $0.536$ & $0.900$ & $1.05$ \\
two-plateau         & $0.003$ & $0.103$ & $0.720$ & $0.958$ & $2.64$ \\
gap                 & $0.003$ & $0.103$ & $0.541$ & $0.910$ & $1.04$ \\
power               & $0.003$ & $0.050$ & $0.533$ & $0.896$ & $1.10$ \\
\bottomrule
\end{tabular}
\end{table}

\paragraph{The multiplicity at large $d$.}  The quadrature of the 13 responses was
measured on 16 cells at $d=64,96$; does it survive the dimensions a real layer has?  It
does: over 96 cells (four families, two $\ell$, three seeds, $d=64,128,256,512$) the ratio
of the quadrature to the measured floor is $0.92$--$1.12$ (median $1.03$) with $97\%$ of
cells inside $\pm10\%$, and at $d=512$ alone $0.997$--$1.080$ (median $1.035$, 24 cells;
Table~\ref{tab:w13} and \texttt{pilot24}).  The responses themselves say why: the mean
$|\cos|$ between the 13 response matrices is $0.017$ at the median and falls with $d$
($0.023$, $0.015$, $0.012$, $0.011$ from $d=64$ to $512$), as an incoherent pair of
$d^2$-dimensional responses should, while the largest single cosine is $0.69$.  The cast carries
$0.85$--$0.86$ of the sum and $C_{\mathrm{total}}$ is flat in $d$ ($3.57$, $3.82$,
$3.67$, $3.80$ for the four dimensions).

\paragraph{The escape.}  The failing runs of \S\ref{sec:escape} repeat one trace.  At
diabetes seed 0, step 24, layer 0 the exact tops are
$0.70,\ 1.87,\ 1.65,\ 1.19,\ 1.01,\ 1.00$ and the E5M2 tops
$0.71,\ 1.93,\ 1.96,\ 3.17,\ 78,\ 1.0\cdot10^{9},\ 4.9\cdot10^{44}$ against certified
upper edges $1.00,\ 1.92,\ 1.71,\ 1.29,\ 1.02,\ 1.00$; the same shape appears at diabetes
seed 2 (step 4) and at digits seeds 0 and 1 (steps 327 and 381).  The step at which a run
escapes is a property of this implementation and this seed --- the same arithmetic escapes
at step 4 for another seed of the same task --- so what repeats is the trace, not its
index.  Beyond a window's edge
the quintic's gain is $2.2\times$ at $1.5$ times the edge and $9.3\times$ at twice, and
the largest start $1+\delta$ whose $p$-bit orbit over the six steps stays below $10$ is
$\delta=0.0075$ at $s=1.01$ for four bits, $0.048$ at $1.035$ and $0.063$ at $1.05$
($0.010$ and $0.035$ at five bits).  With the state rescaled into its window at every step
all 11 runs finish, and all 50 of the long runs below; with the learning rate cut by $10\%$
seven of 11 still fail.  The loop
gain $G=\lambda_{\max}(\mathrm{diag}(\mathrm{step})^{1/2}H\,\mathrm{diag}(\mathrm{step})^{1/2})$,
estimated by power iteration on Hessian-vector products (central differences of the
analytic gradient) every 25 steps, exceeds one on no trajectory (median $\le0.26$; maximum
$0.83$ in five of six; one $7.3$ spike in the sixth).

\section{The Census at 1,200 Steps}
\label{app:census}

\begin{table}[h]
\centering
\caption{The long census (\texttt{pilot25}): 1,200 training steps $\times$ 25 seeds per cell,
the hazard of one Muon call with a 95\% Wilson interval, and the calls at risk.  A call is
one application of the composite to one hidden weight matrix; a run that has not diverged
at 1,200 steps contributes all of its calls and no event (right censoring).  ``clamp'' is
E5M2 at $s=1.01$ with the state rescaled into its window whenever a step would leave it
(\S\ref{sec:escape}).}
\label{tab:census}
\begin{tabular}{llccccc}
\toprule
task & condition & div/runs & calls & hazard & 95\% CI & $1$ in \\
\midrule
\texttt{digits} & E5M2, $s=1.01$   & $25/25$ & $40\,002$  & $6.25\cdot10^{-4}$ & $[4.23,\,9.22]\cdot10^{-4}$ & $1\,600$ \\
        & E5M2, $s=1.035$ & $17/25$ & $99\,822$  & $1.70\cdot10^{-4}$ & $[1.06,\,2.73]\cdot10^{-4}$ & $5\,872$ \\
        & clamp           & $0/25$  & $120\,000$ & $0$ & $[0,\,3.20]\cdot10^{-5}$ & --- \\
\texttt{diabetes} & E5M2, $s=1.01$ & $25/25$ & $2\,730$ & $9.16\cdot10^{-3}$ & $[6.21,\,13.48]\cdot10^{-3}$ & $109$ \\
        & E5M2, $s=1.035$ & $0/25$  & $90\,000$  & $0$ & $[0,\,4.27]\cdot10^{-5}$ & --- \\
        & clamp           & $0/25$  & $90\,000$  & $0$ & $[0,\,4.27]\cdot10^{-5}$ & --- \\
\bottomrule
\end{tabular}
\end{table}

\paragraph{What the run count cannot say.}  At this horizon the two tasks are
indistinguishable as runs --- 25 of 25 fail in both at $s=1.01$, so Fisher's exact test is
undefined --- and they are $14.7\times$ apart per call: $42\,732$ calls carry 50 events, and
the exact conditional binomial test on that split gives $p=3.5\cdot10^{-17}$.  The task is
not a single scale factor either: at $1.035$ \texttt{digits} still fails at $1$ in
$5\,872$ calls (CI $1$ in $9\,400$ to $1$ in $3\,670$) while \texttt{diabetes} shows no
failure in $90\,000$ calls (upper bound $1$ in $23\,400$), so the ordering of the two tasks
is not preserved by the safety --- we report that as measured and offer no mechanism for
it.  What the safety and the clamp do is bounded the same way: $1.035$ lowers the
\texttt{digits} hazard $3.7\times$ and \texttt{diabetes} below $1$ in $23\,400$, the clamp
below $1$ in $31\,200$ on \texttt{digits} and $1$ in $23\,400$ on \texttt{diabetes} over
$210\,000$ calls in all.

\paragraph{The shape of the hazard in time.}  It is not stationary.  On \texttt{digits} at
$s=1.01$ the events per 200-step block are $3$, $10$, $11$, $0$, $1$ over $19\,273$,
$15\,740$, $4\,060$, $800$ and $129$ calls ($0.016\%$, $0.064\%$, $0.271\%$, $0.775\%$ in
the blocks that still have runs at risk), and at $1.035$ the first $400$ steps carry none
of the $17$ events.  On \texttt{diabetes} every event is in the first 200 steps (median
$35$ steps, mean $36.8$): that task fails early or not at all, which is why its $400$-step
census already caught it and why its $1.035$ row is empty.  A run that has survived the
first few hundred steps is therefore not safe, only slower --- the reading a per-run count
of $1,200$ steps cannot give and a per-call rate can.

\paragraph{The prefix check.}  Every 1,200-step run reproduces the $400$-step census of
\texttt{pilot23} on its own prefix: 24 of 24 comparable runs agree (five \texttt{digits}
seeds and three \texttt{diabetes} seeds, three conditions each), which is what makes the
two censuses one experiment and not two.  The script writes no wall-clock field, so the
JSON is a function of the seeds and reruns byte for byte.

\paragraph{Two caveats on the estimate.}  The Wilson interval treats each call as an
independent Bernoulli trial; within a run the calls are correlated through the weights, so
the interval describes the exposure rather than upper-bounding a dependent process.  The
clamp is not free: the implementation costs a four-iteration power estimate per step, which
this paper does not price --- doing so at scale is the one measurement we would add to it.

\section{Additional Tables}
\label{app:tables}

Table~\ref{tab:firstorder} is the square-grid test of
Theorem~\ref{thm:firstorder}; Table~\ref{tab:twochannel} is the two-operand precision
sweep of \S\ref{sec:alloc}; Tables~\ref{tab:variants}--\ref{tab:ci} are the kernel
comparison of \S\ref{sec:variants}, the scale laws, the crossing rule, the two-operand
step-count comparison, the graded schedule against the premium, the fp8 comparison of
\S\ref{sec:fp8}, and the bootstrap intervals; Table~\ref{tab:dsinv} normalises the floor constant by
$\sqrt d$, and Table~\ref{tab:e5m2} is the four-bit repair sweep of \S\ref{sec:fp8}.
Table~\ref{tab:profile} is the per-slot gain profile behind \S\ref{sec:alloc},
Table~\ref{tab:ell} the wrong-dynamic-range decisions behind \S\ref{sec:ellsens},
Table~\ref{tab:realtrain} the seven training conditions of \S\ref{sec:realtrain},
Table~\ref{tab:steps} the step budget of \S\ref{sec:steps}, and
Table~\ref{tab:rules} collects the four rules of \S\ref{sec:rules}.

\paragraph{Scale laws.}  Table~\ref{tab:scaling} gives $S_{\mathrm{peak}}(\ell)$ and the
plateau constant; beside them, measured $\norm{\delta X}_2/u = 1.26$--$1.86$ and
$\norm{\delta B}_2/u = 0.96$--$5.5$ carry no $\sqrt d$ component, while $\mathrm{err}/u$
grows $1:1.07:1.16$ from $d=64$ to $128$ against $1:1.41:2.00$ for a bare $\sqrt d$ law
--- the $\sqrt d$ appears only inside the normalisation of the pair weight
(\S\ref{sec:loop}).

\begin{table}[h]
\centering
\caption{Per-slot gain profile $\beta_t$ of the subspace channel, $d=64$, $p=8$,
$K=6$.  The profile is independent of $K$: the input gain is $11.5$ and $59.3$
at $\ell=10^{-2},10^{-3}$ for every $K\in\{4,6,8,10\}$ tested, and the tail sits at
$1.1u$.}
\label{tab:profile}
\begin{tabular}{rccccccc}
\toprule
$\ell$ & input & step 0 & step 1 & step 2 & step 3 & step 4 & step 5 \\
\midrule
$10^{-2}$ & \textbf{11.5} & 6.2 & 3.1 & 1.6 & 1.1 & 1.1 & 1.1 \\
$10^{-3}$ & \textbf{59.3} & 29.0 & 9.8 & 6.4 & 2.1 & 1.3 & 1.1 \\
$10^{-4}$ & \textbf{512} & 222 & 55.5 & 18.0 & 8.7 & 4.1 & 1.9 \\
\bottomrule
\end{tabular}
\end{table}

\begin{table}[h]
\centering
\caption{The stopping rule with a wrong dynamic range, $\hat\ell=\alpha\ell$: 36 cells
and three targets each, the targets set relative to the true predicted floor.
``missed'' is a returned step count whose measured error exceeds $1.5\varepsilon$;
``refused'' counts cells where the rule refuses a target that some step count of the
true design meets within $1.5\varepsilon$ (each row also holds two cells where no step
count does).  ``understatement'' is $\max\mathrm{err}/J(K)$ on the believed window: how
far the caller's own prediction falls short of the measurement.}
\label{tab:ell}
\begin{tabular}{rcccccc}
\toprule
$\alpha$ & met & missed & refused & $\mathrm{err}/\varepsilon$ median & max & understatement \\
\midrule
$0.1$  & 18 & 0  & 16 & 0.30 & 0.99 & 1.0 \\
$0.25$ & 18 & 0  & 16 & 0.53 & 0.99 & 1.0 \\
$0.5$  & 26 & 0  & 8  & 0.32 & 0.98 & 1.0 \\
$1$    & 36 & 0  & 0  & 0.59 & 1.09 & 1.3 \\
$2$    & 20 & 16 & 0  & 1.25 & 59   & 101 \\
$4$    & 14 & 22 & 0  & 2.52 & 121  & 614 \\
$10$   & 12 & 24 & 0  & 3.11 & 244  & 1275 \\
\bottomrule
\end{tabular}
\end{table}

\begin{table}[h]
\centering
\caption{Theorem~\ref{thm:firstorder} against measurement at the input cast.
$C = \mathrm{sub}_{\mathrm{err}}/u$ measured; $C_{\mathrm{pred}}$ from
\eqref{eq:firstorder} with $E = q(\hat M,p)-\hat M$.  The cell marked $\dagger$
saturates at $\mathrm{sub}_{\mathrm{err}} = 2$.}
\label{tab:firstorder}
\begin{tabular}{rrrccc}
\toprule
$d$ & $\ell$ & $p$ & $C$ measured & $C_{\mathrm{pred}}$ & ratio \\
\midrule
32 & $10^{-1}$ & 8 & 2.24 & 2.24 & 1.00 \\
32 & $10^{-3}$ & 8 & 58.96 & 57.92 & 0.98 \\
32 & $3\cdot10^{-4}$ & 8 & $512.0^{\dagger}$ & 181.5 & $0.35^{\dagger}$ \\
64 & $10^{-1}$ & 8 & 2.14 & 2.14 & 1.00 \\
64 & $10^{-2}$ & 8 & 11.50 & 11.48 & 1.00 \\
64 & $10^{-3}$ & 8 & 59.30 & 58.04 & 0.98 \\
64 & $10^{-3}$ & 16 & 76.28 & 76.26 & 1.00 \\
64 & $3\cdot10^{-4}$ & 11 & 230.25 & 225.74 & 0.98 \\
128 & $10^{-3}$ & 8 & 77.51 & 77.05 & 0.99 \\
128 & $3\cdot10^{-4}$ & 8 & 207.79 & 238.66 & 1.15 \\
\bottomrule
\end{tabular}
\end{table}

\begin{table}[h]
\centering
\caption{Error $\|X_K - UV^\top\|_2$ against the storage precision of the state
$p_X$ (rows) and of the intermediate $p_B$ (columns); shipping design, $\ell=10^{-3}$,
$d=64$, $K=6$.  Each row and column is flat over a wide range and then falls
steeply --- the signature of a bottleneck that moves.}
\label{tab:twochannel}
\begin{tabular}{r|ccccccc}
\toprule
$p_X \backslash p_B$ & 6 & 8 & 10 & 11 & 13 & 16 & 24 \\
\midrule
6  & \e{1.52}{0} & \e{1.45}{0} & \e{1.39}{0} & \e{1.41}{0} & \e{1.40}{0} & \e{1.41}{0} & \e{1.40}{0} \\
8  & \e{3.53}{-1} & \e{3.38}{-1} & \e{3.07}{-1} & \e{3.00}{-1} & \e{3.19}{-1} & \e{3.17}{-1} & \e{3.15}{-1} \\
10 & \e{1.78}{-1} & \e{9.48}{-2} & \e{9.58}{-2} & \e{8.96}{-2} & \e{8.91}{-2} & \e{9.15}{-2} & \e{9.21}{-2} \\
11 & \e{1.55}{-1} & \e{4.26}{-2} & \e{3.48}{-2} & \e{3.73}{-2} & \e{3.43}{-2} & \e{3.49}{-2} & \e{3.51}{-2} \\
13 & \e{1.57}{-1} & \e{2.97}{-2} & \e{1.13}{-2} & \e{1.07}{-2} & \e{1.09}{-2} & \e{9.96}{-3} & \e{9.77}{-3} \\
16 & \e{1.52}{-1} & \e{2.87}{-2} & \e{8.33}{-3} & \e{4.83}{-3} & \e{2.37}{-3} & \e{2.14}{-3} & \e{2.12}{-3} \\
24 & \e{1.56}{-1} & \e{2.90}{-2} & \e{8.62}{-3} & \e{4.81}{-3} & \e{1.90}{-3} & \e{1.64}{-3} & \e{1.62}{-3} \\
\bottomrule
\end{tabular}
\end{table}

\begin{table}[h]
\centering
\caption{Predicted/measured error ratio of the shipped design under four kernels of
$p$-bit arithmetic (\S\ref{sec:variants}), at $K=6$, one matrix per
$(\ell,d)$.  ``eager'' rounds every matrix product, ``fp32 acc.'' accumulates in
float32, ``stochastic'' uses stochastic rounding (mean of three seeds).}
\label{tab:variants}
\begin{tabular}{rrrcccc}
\toprule
$\ell$ & $d$ & $p$ & kernel & eager & fp32 acc. & stochastic \\
\midrule
$10^{-2}$ & 64 & 8 & 0.83 & 0.45 & 0.84 & 0.62 \\
$10^{-2}$ & 64 & 11 & 0.94 & 0.47 & 0.92 & 0.61 \\
$10^{-2}$ & 64 & 16 & 0.82 & 0.50 & 0.86 & 0.56 \\
$10^{-3}$ & 64 & 8 & 1.29 & 0.65 & 1.29 & 0.75 \\
$10^{-3}$ & 64 & 11 & 1.29 & 0.76 & 1.29 & 0.82 \\
$10^{-3}$ & 64 & 16 & 1.28 & 1.04 & 1.24 & 1.21 \\
$10^{-3}$ & 128 & 8 & 0.98 & 0.58 & 0.98 & 0.65 \\
$10^{-3}$ & 128 & 11 & 0.94 & 0.64 & 0.93 & 0.72 \\
$10^{-3}$ & 128 & 16 & 1.33 & 0.93 & 1.32 & 1.07 \\
$10^{-4}$ & 64 & 8 & 1.70 & 1.68 & 1.70 & 1.70 \\
$10^{-4}$ & 64 & 11 & 1.45 & 1.14 & 1.43 & 1.12 \\
$10^{-4}$ & 64 & 16 & 1.03 & 1.02 & 1.03 & 1.02 \\
\bottomrule
\end{tabular}
\end{table}

\begin{table}[h]
\centering
\caption{Graded schedule against uniform precision at equal total bits, as a function
of the first-slot premium.  First row: the real momentum sample; rows 2--5: a
logspace spectrum of the same $128\times32$ shape, two seeds, $K\in\{6,8\}$.  At
$\ell=10^{-4}$ the design's own exact error $E_K$ dominates at the two smaller budgets,
so the gain collapses there and returns at the largest.}
\label{tab:premium}
\begin{tabular}{lcccc}
\toprule
spectrum & premium (bits) & $K{=}6$, 8 bits & $K{=}6$, 12 bits & $K{=}8$, 12 bits \\
\midrule
real $128\times32$ & 0.9 & 1.00 & 1.00 & 1.00 \\
logspace $10^{-1}$ & 1.6--1.7 & 1.09--1.19 & 1.07--1.16 & 1.05--1.13 \\
logspace $10^{-2}$ & 4.3--4.5 & 3.17--3.71 & 3.22--3.59 & 4.23--4.86 \\
logspace $10^{-3}$ & 7.3--7.5 & 6.38--8.10 & 9.13--9.79 & 10.2--15.4 \\
logspace $10^{-4}$ & 8.8--9.8 & 2.25--2.32 & 1.30--1.46 & 20.7--29.1 \\
\bottomrule
\end{tabular}
\end{table}

% scale laws table moved here from the body
\begin{table}[h]
\centering
\caption{Scale laws of the certificate and the floor (synthetic logspace spectra,
shipped design, $p=10$).  $S_{\mathrm{peak}}$ is the maximum of $S(K)$ over $K$,
with the local log--log slope against $\ell$; $C = \mathrm{err}/u$ is the plateau
constant, blank where the design has not converged by $K=8$.}
\label{tab:scaling}
\begin{tabular}{rcccccc}
\toprule
$\ell$ & $S_{\mathrm{peak}}$ & $K^\star$ & local slope & $C$ ($d=32$) & $C$ ($d=64$) & $C$ ($d=128$) \\
\midrule
$10^{-2}$ & 77.1 & 3 & $-1.67$ & 10.1 & 12.2 & 13.4 \\
$10^{-3}$ & 609 & 5 & $-1.88$ & 48.5 & 82.9 & 93.1 \\
$10^{-4}$ & $5.63\times10^{3}$ & 6 & $-1.81$ & -- & -- & -- \\
\bottomrule
\end{tabular}
\end{table}

\begin{table}[h]
\centering
\caption{The crossing rule \eqref{eq:cross} and the target rule \eqref{eq:target} on
frozen real momentum spectra (500 matrices, design window $\ell=10^{-3}$; $K$ swept to
10).  Costs are the measured error at the selected $K$ over the smallest measured error
in $K$.  At $p=24$ the crossing lands on the argmin while the target rule stops three
steps later, at the plateau.}
\label{tab:cross}
\begin{tabular}{rrc rrrc rr}
\toprule
$p$ & plateau (u) & $K_{\arg\min}$ & $K_{\mathrm{cross}}$ & cost & $K_{\mathrm{target}}$ & cost & $K_{\mathrm{target}}$, calibrated $c$ & cost \\
\midrule
6  & 15.8 & 6 & 2 & $4.01\times$ & 10 & $1.002\times$ & 6 & $1.000\times$ \\
8  & 16.0 & 6 & 4 & $9.11\times$ & 10 & $1.001\times$ & 6 & $1.000\times$ \\
11 & 16.0 & 7 & 5 & $16.7\times$ & 10 & $1.001\times$ & 6 & $1.006\times$ \\
16 & 16.0 & 7 & 6 & $6.68\times$ & 10 & $1.001\times$ & 7 & $1.000\times$ \\
24 & 16.1 & 7 & 7 & $1.00\times$ & 10 & $1.006\times$ & 7 & $1.000\times$ \\
\bottomrule
\end{tabular}
\end{table}

\begin{table}[h]
\centering
\caption{Equal-total-bit comparison of the two operands as a function of the step count.
Rows 1--2 hold the design and the matrix fixed and change only $K$; rows 3--4 are the
shipped quintic at $\ell=0.3$, $d=64$; the last row is the frozen real-spectrum replay at
$K=8$.  ``asymmetry'' is the ratio of the two off-diagonal entries.}
\label{tab:k1}
\footnotesize
\setlength{\tabcolsep}{3pt}
\begin{tabular}{>{\raggedright\arraybackslash}p{3.25cm}
                >{\raggedright\arraybackslash}p{2.0cm}
                ccccc}
\toprule
setting & $\ell$, $d$ & balanced & intermediate richer & state richer & ratios & asymmetry \\
\midrule
degree-17, $K=1$, 16 bits & 0.3, 64 & $(8,8)$: \e{1.13}{-2} & $(6,10)$: \e{3.29}{-2} & $(10,6)$: \e{3.33}{-2} & 2.90, 2.94 & $1.01\times$ \\
degree-17, $K=2$, 16 bits & 0.3, 64 & $(8,8)$: \e{1.03}{-2} & $(6,10)$: \e{3.38}{-2} & $(10,6)$: \e{8.32}{-2} & 3.29, 8.11 & $2.46\times$ \\
quintic, $K=2$, 16 bits & 0.3, 64 & $(8,8)$: \e{1.09}{-2} & $(6,10)$: \e{3.79}{-2} & $(10,6)$: \e{1.67}{-2} & 3.48, 1.53 & $2.27\times$ \\
quintic, $K=6$, 16 bits & 0.3, 64 & $(8,8)$: \e{1.21}{-2} & $(6,10)$: \e{4.04}{-2} & $(10,6)$: \e{1.78}{-2} & 3.33, 1.47 & $2.27\times$ \\
real spectra, $K=8$, 32 bits & 0.15, $128\times32$ & $(16,16)$: \e{2.44}{-4} & $(8,24)$: \e{2.82}{-2} & $(24,8)$: \e{5.74}{-2} & 116, 236 & $2.03\times$ \\
\bottomrule
\end{tabular}
\end{table}

\begin{table}[h]
\centering
\caption{Critical safety $s^\star$: the smallest damping at which no run of the
12-spectrum census diverges ($d\in\{48,64\}$, six seeds, $K=6$).  A dash marks a cell
that is already stable at $s=1.0$; the last column is the measured
$c=(s^\star-1)/2^{-p}$.}
\label{tab:safety}
\begin{tabular}{rccc c}
\toprule
$p$ & $\ell=10^{-1}$ & $\ell=10^{-2}$ & $\ell=10^{-3}$ & $c$ \\
\midrule
4  & 1.0028 & 1.0164 & 1.0313 & 0.04--0.50 \\
5  & -- & 1.0053 & 1.0060 & 0.17--0.19 \\
6  & -- & 1.0032 & 1.0047 & 0.20--0.30 \\
8  & -- & 1.0002 & 1.0005 & 0.05--0.12 \\
10 & -- & -- & 1.0001 & 0.08 \\
12 & -- & -- & 1.00001 & 0.04 \\
16 & -- & -- & -- & -- \\
\bottomrule
\end{tabular}
\end{table}

\begin{table}[h]
\centering
\caption{The floor constant with the pair operator normalised by $\sqrt d$
(\S\ref{sec:loop}): measured $(\mathrm{err}-E_K)/(u\,\hat v\norm{W}_F/\sqrt d)$ with
$W=[1/(\sigma_i+\sigma_j)]$ and $\hat v=\norm{\hat M}_F/(\sqrt3 d)$, median over two
$\ell$, two $p$ and three seeds per cell.  The last column is the held-out grid
($\ell\in\{5\cdot10^{-2},3\cdot10^{-3}\}$, $p\in\{6,16\}$, other seeds, $d\in\{48,128\}$).
The two columns after $d=256$ extend the same grid to four times its top, and at
$d=768$ and $1024$ the two flat-start families read $7.94$/$7.93$ and $6.11$/$6.17$.
Without the $\sqrt d$ these ratios fall with exponent $-0.42$ in $d$; weighting the
operator by $\sigma_i\sigma_j/(\sigma_i+\sigma_j)$ instead spreads $65\times$
($2.6\times$ worse out of sample) and normalising by its stable rank $136\times$
($2.9\times$).}
\label{tab:dsinv}
\begin{tabular}{lcccccccc}
\toprule
family & $d=32$ & $d=48$ & $d=64$ & $d=128$ & $d=256$ & $d=384$ & $d=512$ & held out \\
\midrule
logspace & 6.64 & 6.77 & 6.96 & 7.33 & 7.64 & 7.76 & 7.75 & 6.96 \\
two-plateau & 4.85 & 5.14 & 5.63 & 5.89 & 5.96 & 6.04 & 6.13 & 5.37 \\
gap & 5.44 & 5.56 & 5.85 & 6.11 & 6.27 & 6.40 & 6.40 & 5.96 \\
power & 5.30 & 6.00 & 6.50 & 7.35 & 7.25 & 7.49 & 7.74 & 6.26 \\
\bottomrule
\end{tabular}
\end{table}

\begin{table}[h]
\centering
\caption{Four-bit repairs at $K=6$ with the shipped design, 12-spectrum census
($d\in\{48,64\}$, six seeds).  ``blow-up'' counts runs whose state distance
$\norm{X_K-UV^\top}_2$ exceeds $5$; ``floor'' is the median subspace error of the census
(cap $2$); $s^\star$ is the smallest damping with no blow-up.  Each row names the
precision of the slots in order; ``cast'' is the input quantisation, E4M3 is $p=5$ and
bf16 $p=8$.}
\label{tab:e5m2}
\begin{tabular}{lcccccc}
\toprule
rule & blow-up & floor & floor & blow-up & floor & $s^\star$ \\
     & $\ell=10^{-2}$ & $10^{-2}$ & $10^{-1}$ & $\ell=10^{-3}$ & $10^{-3}$ & $10^{-3}$ \\
\midrule
$p=4$, $s=1.01$ & 4/12 & 0.877 & 0.217 & 12/12 & 2.00 & 1.0313 \\
$p=4$, $s=1.035$ & 0/12 & 0.841 & 0.221 & 0/12 & 2.00 & 1.0313 \\
cast bf16 & 5/12 & 0.419 & 0.165 & 10/12 & 2.00 & 1.0235 \\
cast, step 1 bf16 & 0/12 & 0.243 & 0.133 & 10/12 & 2.00 & 1.0221 \\
cast, steps 1--2 bf16 & 0/12 & 0.146 & 0.114 & 5/12 & 0.537 & 1.0238 \\
cast, step 1 E4M3 & 1/12 & 0.469 & 0.157 & 9/12 & 2.00 & 1.0219 \\
cast, steps 1--2 E4M3 & 0/12 & 0.416 & 0.140 & 9/12 & 2.00 & 1.0221 \\
E4M3 throughout & 0/12 & 0.411 & 0.107 & 0/12 & 1.92 & 1.0060 \\
bf16 throughout & 0/12 & 0.049 & 0.014 & 0/12 & 0.283 & -- \\
\bottomrule
\end{tabular}
\end{table}

\paragraph{The (format, step-count) grid.}  On \texttt{diabetes} (three seeds, 1200
steps) the training MSE at K=2 is $1.71\times10^{-3}$ for bf16 (24 bits, update error
$0.71$), $1.46\times10^{-3}$ for E4M3 (15 bits, $0.71$) and $2.14\times10^{-3}$ for E5M2
at safety $1.035$ (12 bits, $0.73$); at K=4 they are $1.58$, $1.60$ and
$1.62\times10^{-3}$, and no cell at $K\le4$ diverges.  Against the fp32 value $1.299\times10^{-3}$ the spread over the whole grid is
a factor $1.65$, while the update error spans two orders of magnitude.

\begin{table}[h]
\centering
\caption{The kernel inside training (\S\ref{sec:realtrain}).  Update error is the window
error at the last probe, median over seeds; ``div'' counts weight blow-up
($\max|W|>10^{3}$, reached before NaN).  Bits are the cast plus one stored width per
step.  All rows are $400$ steps; the $1{,}200$-step census is Table~\ref{tab:census}.
The E5M2 row at safety $1.01$ has no diabetes loss: all three runs stop.}
\label{tab:realtrain}
\begin{tabular}{lccccccc}
\toprule
condition & bits & upd.\ err & \texttt{digits} train & acc & div & \texttt{diabetes} train & div \\
\midrule
fp32              & 168 & $2.8\cdot10^{-14}$ & $5.84\cdot10^{-5}$ & $0.980$ & $0/5$ & $1.30\cdot10^{-3}$ & $0/3$ \\
bf16              & 56  & $1.6\cdot10^{-2}$  & $5.79\cdot10^{-5}$ & $0.977$ & $0/5$ & $1.38\cdot10^{-3}$ & $0/3$ \\
E4M3              & 35  & $1.2\cdot10^{-1}$  & $6.07\cdot10^{-5}$ & $0.980$ & $0/5$ & $1.46\cdot10^{-3}$ & $1/3$ \\
E5M2, $s=1.01$    & 28  & $2.4\cdot10^{-1}$  & $7.39\cdot10^{-5}$ & $0.977$ & $2/5$ & --- & $3/3$ \\
E5M2, $s=1.035$   & 28  & $2.5\cdot10^{-1}$  & $6.22\cdot10^{-5}$ & $0.980$ & $0/5$ & $1.57\cdot10^{-3}$ & $0/3$ \\
hybrid, E5M2 body & 40  & $1.4\cdot10^{-1}$  & $5.77\cdot10^{-5}$ & $0.980$ & $0/5$ & $1.44\cdot10^{-3}$ & $0/3$ \\
hybrid, E4M3 body & 44  & $7.3\cdot10^{-2}$  & $5.83\cdot10^{-5}$ & $0.977$ & $0/5$ & $1.42\cdot10^{-3}$ & $0/3$ \\
\bottomrule
\end{tabular}
\end{table}

\begin{table}[h]
\centering
\caption{Step count against precision on 80 real momentum matrices (\S\ref{sec:steps}),
design window $\ell_0=10^{-2}$; bits are the input cast plus one stored width per step,
and the error is the window error, median over the 80 matrices.  The last two rows are the
hybrid rule of \S\ref{sec:fp8} at $K=6$.}
\label{tab:steps}
\begin{tabular}{llccc}
\toprule
format & $K$ & bits & window error & \\
\midrule
fp32    & 6 & 168 & $2.8\cdot10^{-14}$ & \\
\midrule
E5M2    & 2 & 12 & $7.35\cdot10^{-1}$ & \\
        & 3 & 16 & $3.64\cdot10^{-1}$ & \\
        & 4 & 20 & $2.49\cdot10^{-1}$ & \\
        & 5 & 24 & $2.41\cdot10^{-1}$ & plateau \\
        & 6 & 28 & $2.38\cdot10^{-1}$ & \\
\midrule
E4M3    & 3 & 20 & $3.13\cdot10^{-1}$ & \\
        & 4 & 25 & $1.29\cdot10^{-1}$ & knee \\
        & 5 & 30 & $1.24\cdot10^{-1}$ & plateau \\
        & 6 & 35 & $1.25\cdot10^{-1}$ & \\
\midrule
bf16    & 2 & 24 & $7.17\cdot10^{-1}$ & \\
        & 3 & 32 & $2.96\cdot10^{-1}$ & \\
        & 4 & 40 & $2.43\cdot10^{-2}$ & knee \\
        & 5 & 48 & $1.60\cdot10^{-2}$ & plateau \\
        & 6 & 56 & $1.62\cdot10^{-2}$ & \\
\midrule
hybrid, E5M2 body & 6 & 40 & $1.45\cdot10^{-1}$ & \\
hybrid, E4M3 body & 6 & 44 & $7.36\cdot10^{-2}$ & \\
\bottomrule
\end{tabular}
\end{table}

\begin{table}[h]
\centering
\caption{The four rules of \S\ref{sec:rules}, the measurement behind each and its price.
``bits'' is the width of the input cast plus one stored width per step; a named dynamic
range means the logspace ladder at $d=64$, four seeds.}
\label{tab:rules}
\begin{tabular}{p{3.2cm}p{4.6cm}p{2.5cm}p{3.2cm}}
\toprule
decision & measurement & window error & cost \\
\midrule
stop at the plateau ($K=5$) & flat from $K=5$ in every format; sixth step $<1.5\%$; $78/80$ matrices prefer it & $1.60\cdot10^{-2}$ (bf16) & $+8$ bits, no gain \\
buy steps, not precision & $40$-bit hybrid $1.45\cdot10^{-1}$ against bf16 $K=4$ $2.43\cdot10^{-2}$; one downgraded slot $2.7$--$12.6\times$ & $6\times$ at equal bits & $4$ bits per slot saved \\
$p=4$ only with $s\ge1.035$ & blow-up $12/12\to0/12$ at $10^{-3}$; real data $6/11\to0/11$ at $400$ steps, $17/25$ at $1{,}200$ & cap $2.00$ in both & $+1.9\%$ at $10^{-1}$, $-4.1\%$ at $10^{-2}$ \\
$\hat\ell\le\ell$ (halve $\sigma_{\min}$) & error $\le1.07\times$, $\le1$ extra step & --- & one extra step \\
\bottomrule
\end{tabular}
\end{table}

\begin{table}[h]
\centering
\caption{fp8 against the same kernel with an unlimited exponent at the matched precision
(\S\ref{sec:fp8}; median of five seeds, $d=64$, $K=6$).  ``E4M3 scaled'' uses
per-tensor power-of-two scales; ``sub'' counts the stored singular directions of the
start that fall below the format's smallest normal.  $\dagger$ marks a cell whose error
sits at the cap $2$ that any pair of orthogonal factors obeys; at $\ell=1.3\times10^{-2}$
E5M2 saturates $8192$ stored values.}
\label{tab:fp8}
\begin{tabular}{rrccccc}
\toprule
$\ell$ & $p=5$ ideal & E4M3 & E4M3 scaled & $p=4$ ideal & E5M2 & sub \\
\midrule
$0.15$   & \e{9.86}{-2} & \e{9.89}{-2} & \e{1.03}{-1} & \e{1.85}{-1} & \e{1.84}{-1} & 0 \\
$0.1$    & \e{1.08}{-1} & \e{1.11}{-1} & \e{1.14}{-1} & \e{2.15}{-1} & \e{2.15}{-1} & 0 \\
$3.6\cdot10^{-2}$ & \e{1.69}{-1} & \e{1.71}{-1} & \e{1.71}{-1} & \e{3.48}{-1} & \e{3.58}{-1} & 0 \\
$1.3\cdot10^{-2}$ & \e{3.00}{-1} & \e{3.08}{-1} & \e{8.08}{-1} & $2.00^{\dagger}$ & $2.00^{\dagger}$ & 5 \\
$4.8\cdot10^{-3}$ & \e{7.21}{-1} & \e{7.67}{-1} & \e{1.40}{0} & $2.00^{\dagger}$ & $2.00^{\dagger}$ & 14 \\
\bottomrule
\end{tabular}
\end{table}

\begin{table}[h]
\centering
\caption{Bootstrap 95\% percentile intervals (10\,000 resamples, fixed seed) on every
ratio and constant the abstract and the introduction quote.  ``point'' is the statistic on the
observed sample.  The 500-cell intervals are narrow because the formula reproduces that
grid to three digits; the interval says the agreement is not carried by a few cells.}
\label{tab:ci}
\begin{tabular}{llcc}
\toprule
quantity & cells & point & interval \\
\midrule
certificate, err/pred median & 192 & 0.99 & $[0.98,1.00]$ \\
certificate, worst deviation & 192 & $1.60$ & $[1.59,1.60]$ \\
stopping rule, pred/meas median ($d=64$) & 90 & $1.11$ & $[1.07,1.14]$ \\
stopping rule, pred/meas median ($d=128$) & 90 & $1.10$ & $[1.08,1.13]$ \\
two-term formula, pred/meas median ($p=8$) & 500 & $1.0000$ & $[0.9999,1.0001]$ \\
self-calibrated prediction, held out & 9 & $0.85$ & $[0.64,1.12]$ \\
graded schedule, gain over uniform & 48 & $3.06$ & $[2.24,4.76]$ \\
certificate on held-out shapes and spectra & 48 & $0.42$ & $[0.37,0.48]$ \\
first-order formula on the same cells & 48 & $0.89$ & $[0.81,0.91]$ \\
dimension-normalised constant, fit & 240 & $6.07$ & $[5.98,6.14]$ \\
the same, held-out cells & 64 & $6.04$ & $[5.78,6.26]$ \\
\bottomrule
\end{tabular}
\end{table}

\paragraph{Discretisation checks.}
Two of the paper's quantities are produced by linear programs and one by integer
rounding, so each could be an artefact of its grid.  (i) The minimax single-step error at
$\ell=10^{-3}$ moves by $7.5\times10^{-4}$ (degree 5) and $2.2\times10^{-3}$ (degree 9)
between $n=50$ and $n=3200$ grid points, and by $2\times10^{-5}$ between $n=1600$ and
$n=3200$; degree 17 needs $n\ge400$ (a $5\%$ deviation at $n=50$, $1.9\times10^{-5}$ at
$n=3200$), and the values quoted in the tables use $n\ge400$.  (ii) The last-step family's coefficients at $p'(1)=\lambda$ move by
$5\times10^{-5}$ between $n=100$ and $n=1600$, and the argmin of $J$ over $\lambda$ is the
same grid point on $n\in\{100,400,1600\}$ with $J$ changing by $4\times10^{-6}$.
(iii) Against the exact integer optimum from a dynamic program over the $K+1$ slots, the
largest-remainder schedule of \eqref{eq:maxform} is exactly optimal in 44 of the 48
cells and within $1.47\times$ in the rest, so the quoted gain is not a rounding artefact.

\paragraph{A verification in fp8.}
Table~\ref{tab:fp8} also documents where the format leaves the law's domain: from
$\ell=1.3\times10^{-2}$ downward the smallest stored directions are subnormal (5 of 64
there, 14 at $4.8\times10^{-3}$), and below $\ell\approx6\times10^{-4}$ directions are
flushed and the error sits at the cap in every format, including the unlimited-exponent
one, because four significant bits are not enough for a dynamic range of $10^{3}$.
% --- END sections/99-appendix.tex ---


\bibliographystyle{plain}
\bibliography{references}

\end{document}
